\documentclass[11pt]{article}
\usepackage[T1]{fontenc}
\usepackage[utf8]{inputenc}
\usepackage{lmodern,amsmath,amssymb,amsthm,mathtools,booktabs,array}
\usepackage[margin=1in]{geometry}
\usepackage{xurl}
\usepackage[hidelinks]{hyperref}
\usepackage{microtype}
\hypersetup{pdftitle={Finite games that read the type of a von Neumann algebra},pdfauthor={Seth Douglas and Nidhal Mghirbi}}
\newcommand{\DoNotLoadEpstopdf}{}
\newcommand{\R}{\mathbb R}
\newcommand{\N}{\mathbb N}
\newcommand{\Z}{\mathbb Z}
\newcommand{\C}{\mathbb C}
\newcommand{\Q}{\mathbb Q}
\newcommand{\Tr}{\operatorname{Tr}}
\newcommand{\Sp}{\operatorname{Sp}}
\newcommand{\id}{\operatorname{id}}
\newcommand{\Var}{\operatorname{Var}}
\newcommand{\supp}{\operatorname{supp}}
\newcommand{\FMS}{G_{\mathrm{FMS}}}
\newcommand{\GIII}{G_{\mathrm{III}}}
\newcommand{\Gthree}{G_3}
\newcommand{\eps}{\varepsilon}
\newcommand{\vworst}{v^{\mathrm{worst}}}
\newcommand{\un}{\mathbf1}
\newcommand{\sq}{\sqrt2}
\newtheorem{theorem}{Theorem}[section]
\newtheorem{lemma}[theorem]{Lemma}
\newtheorem{proposition}[theorem]{Proposition}
\newtheorem{corollary}[theorem]{Corollary}
\theoremstyle{definition}
\newtheorem{definition}[theorem]{Definition}
\newtheorem{remark}[theorem]{Remark}
\newtheorem{conjecture}[theorem]{Conjecture}
\setlength{\emergencystretch}{2em}
\raggedbottom
\title{Finite games that read the type\\of a von Neumann algebra\thanks{The authors developed this work with AI models from more than one developer and take responsibility for every claim. How the work was done, and what the machine checks cover, is described at the end of the paper.}}
\author{Seth Douglas \and Nidhal Mghirbi}
\date{Version 1.0.0, October 2026}
\begin{document}
\maketitle
\begin{abstract}
We construct finite two-player games with binary scoring whose normal strategies distinguish operator-algebraic type. A 243-pair game has perfect play in type III and none when either ambient algebra is semifinite. For factors with separable predual, hosting it is equivalent to an invariant normal state on the flow of weights, or to an expected hyperfinite type III1 subfactor. Prepared-state hosting is therefore insensitive to the ratio invariant throughout the injective flow-finite class. A different quantifier recovers that invariant: every faithful normal state has value one for the flagged magic-square game exactly when 1/2 belongs to the Connes invariant; a nine-question conference-basis game gives the corresponding criterion for 1/3. Together these two response properties characterize type III1. One prepared state already wins both games in dimension 25. We prove that every positive rational is forced as a modular ratio by some finite game, and identify the exact rational orthogonal-pair single-bridge parameters. For the type-III-only game, a robust sorted spectral law holds for every normal semifinite strategy. It yields matching entropy, deficit, dimension and variance exponents, including an exponential inverse-square-root dimension law. No validated numerical constant is given; the sharp constant is conjectural. We give explicit finite templates, retain their boundary penalties, and distinguish sorted trace position from log density. The specified normal semifinite perturbative gravitational models cannot attain perfect play, whereas suitable type III quantum-field-theory hosts can. This conditional algebraic consequence gives no preparation or energy bound.
\end{abstract}
\section{Definitions and main results}\label{sec:definitions}
A finite game is a tuple $G=(X,Y,(A_x),(B_y),\pi,V)$ with finite nonempty question and answer sets, a probability distribution $\pi$ on $X\times Y$, and a predicate $V(a,b\mid x,y)\in\{0,1\}$. Its loss is
\[
 \ell_G=\sum_{x,y,a,b}\pi(x,y)(1-V(a,b\mid x,y))p(a,b\mid x,y).
\]
Only pairs of positive question weight matter to perfection. A normal commuting strategy consists of commuting positive operator valued measurements (POVMs) and a state whose marginal on the specified ambient von Neumann algebra is normal. Mixed states and larger commuting partners are allowed.

For a faithful normal state $\phi$ on a $\sigma$-finite von Neumann algebra $M$, use its standard form $(M,L^2(M),J,\mathcal P)$ and natural-cone vector $\xi_\phi$. Alice's effects lie in $M$ and Bob's in $M'=JMJ$. Write
\[
 p(a,b\mid x,y)=\langle\xi_\phi,A_a^xB_b^y\xi_\phi\rangle,
 \qquad v_G(M,\phi)=\sup_{A,B}(1-\ell_G).
\]
\emph{Hosting} means existence of a normal state and a perfect strategy with the standard commutant as partner. A nonfaithful state is interpreted on its support corner. For factors define
\[
 \vworst_G(M)=\inf_{\phi\text{ faithful normal}}v_G(M,\phi).
\]
Thus $\vworst_G(M)=1$ means that for every faithful normal state and every accuracy there are suitable measurements. Measurements may depend on both. It does not assert perfect attainment at that state.

For a faithful state, $M_\phi$ is its centralizer and $\sigma^\phi$ its modular group. An $r$-eigenoperator is a nonzero bounded $x$ with $\sigma_t^\phi(x)=r^{it}x$. Connes's invariant $S(M)$ is the intersection of modular spectra over faithful normal semifinite weights. Its positive part is a closed multiplicative subgroup; for countably decomposable factors the intersection over faithful normal \emph{states} has the same positive part \cite{Connes}. For semifinite factors that part is $\{1\}$; for type $\mathrm{III}_\lambda$, $0<\lambda<1$, it is $\lambda^\Z$; and for type $\mathrm{III}_1$ it is $\R_{>0}$. Type $\mathrm{III}_0$ has $S(M)=\{0,1\}$.

For local dimensions at most $d$ on both sides, $\omega_d(G)$ denotes the supremum winning probability over spatial POVM strategies, including mixed states. Let $D_G(\eps)$ be the smallest such $d$ permitting loss at most $\eps$.

The flow of weights is the dual action $\theta$ on the center of the continuous core $M\rtimes_{\sigma^\phi}\R$. We use the convention $\theta_s(h)=e^{-s}h$ and $\tau\circ\theta_s=e^{-s}\tau$. An expected subfactor is the range of a faithful normal conditional expectation. Write $R_\lambda$ for the injective type $\mathrm{III}_\lambda$ factor and $R_\infty$ for the injective type $\mathrm{III}_1$ factor, with separable predual.

\begin{theorem}[Finite-game hosting]\label{thm:host}
The game $\GIII$ defined below has 45 Alice questions, 46 Bob questions, at most nine answers, and 243 uniformly sampled pairs. It has a perfect normal strategy on $R_{1/2}$. No perfect strategy has a normal marginal on a semifinite ambient algebra of either party. For every factor $M$ with separable predual, the following are equivalent:
\begin{enumerate}
\item $M$ hosts $\GIII$;
\item $M$ contains an expected copy of $R_{1/2}$;
\item its flow of weights has an invariant normal state;
\item $M$ contains an expected copy of $R_\infty$.
\end{enumerate}
The displayed perfect correlation belongs to $C_{qa}$, the closure of finite-dimensional spatial correlations, but has no normal semifinite realization.
\end{theorem}
The game and exclusion are proved in Section~\ref{sec:games}; the expected-subfactor and flow equivalences in Section~\ref{sec:hosting}. The analytic equivalence between expected embeddings and invariant flow states is Haagerup--Musat's theorem, not a new structure theorem here.

\begin{theorem}[Two finite state-response tests]\label{thm:two}
For every factor $M$ with separable predual,
\[
 \vworst_{\FMS}(M)=1\iff 1/2\in S(M),\qquad
 \vworst_{\Gthree}(M)=1\iff 1/3\in S(M).
\]
Consequently every faithful normal state has value one for both games if and only if $M$ is type $\mathrm{III}_1$.
\end{theorem}
Here $\FMS$ has five questions per party, while $\Gthree$ has nine per party and five answers per question. Section~\ref{sec:states} proves the theorem. A faithful state on $M_{25}$ wins both games perfectly. The universal-state premise therefore cannot be replaced by one preparation, even with exact scores.

\begin{theorem}[Rational ratio forcing]\label{thm:ratio-summary}
Every $r\in\Q_{>0}$ is forced by some finite binary-scoring game admitting a finite-dimensional perfect strategy: every perfect normal strategy has a nonzero modular $r$-eigenoperator after support reduction. On the binary magic-square core every rational ray of $\R^4$ is finitely forceable. The distinguished parameters of rational orthogonal-pair single bridges are exactly $\{2+x^2+y^2:x,y\in\Q\}$.
\end{theorem}
Section~\ref{sec:ratios} proves these statements, supplies six explicit games, and states the arithmetic restrictions with their construction hypotheses. Existence of a game forcing each rational does not prescribe every rational Schmidt spectrum or classify all finite-game distinctions.

\begin{theorem}[Robustness and dimension]\label{thm:price}
For a normal semifinite strategy of $\GIII$ with density $h$, let $p(t)=\mu_t(h)$ be its decreasing singular-value function, normalized by $\int p=1$. Put
\[
 H(h)^2=\int_0^\infty\bigl(\sqrt{p(t)}-\sqrt{2p(2t)}\bigr)^2\,dt.
\]
Theorem 7.2 of the companion paper \cite{Base} supplies a universal finite constant $C_F$. There is a universal constant
\[
 C_R=\frac{243}{8}(C_F^2+4552),
\]
such that $H(h)^2\le C_R\ell_{\GIII}$. It holds for either normal semifinite marginal, mixed strategies and POVMs. As $\eps\downarrow0$, the minimum local dimension satisfies
\[
 D_{\GIII}(\eps)=\exp\bigl(\Theta(\eps^{-1/2})\bigr),\qquad
 1-\omega_d(\GIII)=\Theta((\ln d)^{-2}).
\]
The optimized type-I marginal entropy and normalized finite-factor deficit have exponent $1/2$; sorted log trace-position and log-density variances have exponent $1$. The exponents match explicit constructions, not their leading constants.
\end{theorem}
Section~\ref{sec:prices} specifies each resource and proves these statements. No validated numerical constant is given; the sharp constant is conjectural.

\begin{theorem}[Type and dimension cost]\label{thm:phase-summary}
There are finite rational games of minimum host class II with polynomial cost, exponential cost, and a quantum value gap. Every $\Delta^0_2$ exponent in $[1/2,1]$ occurs as an existing log-dimension exponent of a type-III-only game with quantum value one. Every positive finite existing exponent is $\Delta^0_2$. Type-III-only hosting is d.r.e.-hard; its intersection with quantum value one is $\Pi^0_2$-hard as an unpromised language.
\end{theorem}
Section~\ref{sec:phase} proves these claims using published quantitative games, a sharp matrix rank/defect bound and explicit AND/OR constructions. The universal relation between type and dimension cost and the exact arithmetical classification remain open.

\paragraph{Prior work and the contribution.}
The flagged magic-square gadget and frustrated-overlap principle are due to Man\v cinska \cite{FMS}; the conference geometry is classical \cite{Conference}. Four-dimensional Kochen--Specker configurations provide earlier context \cite{Cabello}. Connes's classification, state-orbit theory, and modular expectations provide the analytic tools \cite{Connes,CS,HS,Takesaki}. Exact Cuntz mechanisms and operational readings of type through embezzlement precede this work \cite{Harris,VLSWW}. The contribution is the particular finite-game equivalences, the rational-ray and ratio-forcing theorems, and the robust gluing law for $\GIII$. No earlier version of these precise finite-game statements was found in the bounded comparison of the cited literature; this is not an exhaustive priority claim. The authors' companion work supplies the base game and robust sector estimate \cite{Base}; their other public dimension witnesses already have faster-growing dimension requirements \cite{Explicit,Precision}.

\begin{lemma}[Normal support reduction]\label{lem:support}
Every normal-marginal commuting strategy can be represented on the faithful standard form of its support corner, preserving all correlations.
\end{lemma}
\begin{proof}
Purify the joint state into a commuting extension or use its joint GNS representation. If $p$ is the marginal support, then $p\psi=\psi$ and Bob commutes with $p$, so $A$ can be replaced by $pAp$. On $M=pNp$ the map $x\xi_\phi\mapsto x\psi$ extends to an isometry onto $\overline{M\psi}$. This subspace reduces $M$. Compression of Bob's effects by the isometry gives POVMs in $M'$, even if Bob originally failed to preserve that cyclic subspace. The compressed marginal is faithful and normal. Semifiniteness passes to the corner.
\end{proof}

\section{The predicates and the modular bridge}\label{sec:games}
\subsection{Flagged magic square and its nine-copy gluing}
Let $\mathcal R=\{+++,+--,-+-,--+\}$ and let $\mathcal C_j$ be the triples of product $\eta_j$, where $\eta=(+,+,-)$. Put $r_0=c_0=+++$ and $r_1=c_1=+--$. Alice's questions are $R_1,R_2,R_3,T_0,T_1$; Bob's are $C_1,C_2,C_3,U_0,U_1$. A row answer is in $\mathcal R\cup\{f\}$, and a column answer in $\mathcal C_j\cup\{f\}$. On $T_x$ the answers are $\mathcal C_1\setminus\{c_x\}$ and two signs; on $U_y$ they are $\mathcal R\setminus\{r_y\}$ and two signs.

At $(R_i,C_j)$ accept two flags, reject a single flag, and otherwise require equal intersecting signs. At $(T_x,C_1)$ a sign matches the class $\{f,c_x\}$; a triple must match Bob's triple. At $(R_1,U_y)$ use the reversed rule with $r_y$. At $(T_x,U_y)$ two signs must satisfy $ab=(-1)^{xy}$; if either answer is a triple accept. All other cells accept all valid answers. Sample the 25 pairs uniformly. This is $\FMS$, in a real Pauli convention unitarily equivalent to \cite{FMS}.

For $\GIII$ take nine disjoint copies indexed by $j\in\Z/9\Z$ and one extra Bob question $\mathrm{Lab}$, with answers $a\in\Z/9\Z$. Besides the 225 internal pairs, use
\[
 (R_1^j,\mathrm{Lab}):\quad (\text{Alice answers }f)\iff(a=j),
\]
\[
 (R_2^j,\mathrm{Lab}):\quad (\text{Alice answers }+++)\iff(a\in\{j+1,j+2\}).
\]
These are the 243 equally weighted pairs. Other cross-copy pairs have weight zero. All indices below are modulo nine. The finite predicate is also enumerated in the ancillary data.

The real square used to define the ideal projectors is
\[
 O=\begin{pmatrix}X\otimes I&I\otimes X&X\otimes X\\
 I\otimes Z&Z\otimes I&Z\otimes Z\\
 X\otimes Z&Z\otimes X&Y\otimes Y\end{pmatrix}.
\]
Joint row and column eigenspaces give the four magic outcomes; append a one-dimensional flag. The calibrated rays are
\[
 c_0=(1,0,1,0)/\sq,\quad c_1=(0,1,0,1)/\sq,
\]
\[
 r_0=(1,1,1,1)/2,\quad r_1=(1,-1,1,-1)/2.
\]
Their overlap is $\langle c_x,r_y\rangle=(-1)^{xy}/\sq$. The sign projectors are the rays $(f\pm c_x)/\sq$ and $(f\pm r_y)/\sq$; all remaining outcomes are the calibrated residual magic rays. They win on
\[
 \Psi_2=(|ff\rangle+\sq\sum_{i=0}^3|ii\rangle)/3.
\]

\subsection{The conference-basis game}
For $\Gthree$, both players have base questions $R_1,R_2,R_3,C_1,C_2,C_3,H$. Each has four magic rays and a flag. The six old bases are the joint eigenspaces above. The additional orthonormal basis consists of the columns
\[
 r_0=\frac{(1,1,-1,0)^T}{\sqrt3},\quad
 r_1=\frac{(1,-1,0,1)^T}{\sqrt3},\quad
 r_2=\frac{(0,1,1,1)^T}{\sqrt3},\quad
 r_3=\frac{(1,0,1,-1)^T}{\sqrt3}.
\]
Their numerator Gram matrix is $3I$. Up to signed row and column permutations, these numerators form the classical skew conference matrix of order four, the $q=3$ case of the Paley construction; see Goethals--Seidel \cite[Section 2, Theorem 2.1]{Conference}. In the coordinate basis the $R_2$ labels are $d_0=+++$, $d_1=-+-$, $d_2=+--$, $d_3=--+$. Alice adds $T_0,T_1$, with signs and residual labels $k_j$ for $j\ne x$. Bob adds $U_0,U_1$, with signs and labels $h_j$ for $j\ne y$. Sample all 81 pairs uniformly and use the following rules.
\begin{enumerate}
\item On base/base cells accept two flags and reject one flag. Two magic outcomes are accepted exactly when $\Tr(P_{q,a}P_{r,b})>0$ for the specified rational ray projectors.
\item On $(T_x,R_2)$ a sign matches $\{f,d_x\}$ and $k_j$ matches $d_j$.
\item On $(H,U_y)$ a Bob sign matches $\{f,y\}$ and $h_j$ matches $j$.
\item On $(T_x,U_y)$ two signs obey $ab=(-1)^{xy}$; residual answers are accepted.
\item All remaining cells accept every valid pair.
\end{enumerate}
Thus the predicate uses rational zero tests only. There are five answers at every question and 732 forbidden events. Use $c_0=e_0,c_1=e_1$ and the displayed $r_0,r_1$ for the sign rays. The ideal state is
\[
 \Psi_3=(|ff\rangle+\sqrt3\sum_{i=0}^3|ii\rangle)/\sqrt{13}.
\]
Base orthogonality and calibration give zero forbidden amplitudes. For signs the amplitude is $(1+ab(-1)^{xy})/(2\sqrt{13})$, so the rejected parity amplitudes vanish.

\subsection{Extraction in arbitrary normal strategies}
\begin{lemma}[Checked classes]\label{lem:classes}
A perfectly matched coarse partition is a centralizer PVM, with its unique standard commutant mirror. Fine effects with sharp marginals are the corresponding joint spectral projections.
\end{lemma}
\begin{proof}
Use Lemma~\ref{lem:support}. For different classes $c,d$, the commuting positive product $E_cF_d$ has expectation zero and annihilates $\xi=\xi_\phi$. Summing gives $E_c\xi=E_cF_c\xi=F_c\xi$. Hence $E_cE_d\xi=E_cF_d\xi=0$. Separation and completeness give orthogonality and idempotence. A self-adjoint mirror gives $\phi(E_cx)=\phi(xE_c)$, so $E_c\in M_\phi$ and $F_c=JE_cJ$. A positive fine effect lies under each of its projection marginals; completeness saturates their products.
\end{proof}
Apply this to flag/magic classes and cell signs. On the magic projection $m=1-f$ the square relations hold internally. Write $a=O_{11},b=O_{12},c=O_{21},d=O_{22}$. They give
\[
 [a,b]=[a,c]=[b,d]=[c,d]=0,\qquad cb=-bc,\qquad ad=-da.
\]
Indeed the two expressions for $O_{33}$ are $acbd=-abcd$; multiplication by $a$ and $d$ gives the first anticommutation, and $[ac,bd]=0$ gives the second. Two commuting Pauli pairs therefore generate a unital centralizer $M_4$, with arbitrary multiplicity. No trace or finite-dimensional rigidity assumption has been used.

For $\Gthree$ the diagonal base cells check every outcome. An actual conference outcome $E_j$ is orthogonal to every old ray orthogonal to $r_j$: the forbidden event is $\phi(PE_jP)=0$ and faithfulness gives $E_jP=0$. The following old rays span the four required orthogonal complements:
\[
\begin{array}{c|ccc}
j&\multicolumn{3}{c}{\text{spanning rays}}\\\hline
0&e_3&e_0-e_1&e_0+e_2\\
1&e_2&e_0+e_1&e_0-e_3\\
2&e_0&e_1-e_2&e_1-e_3\\
3&e_1&e_0-e_2&e_0+e_3
\end{array}
\]
Consequently $E_j\le r_jr_j^T\otimes1$. Both four-tuples sum to $m$, so each positive difference is zero. This fixes the conference basis completely, including multiplicities.

\begin{lemma}[Modular bridge]\label{lem:bridge}
Every normal perfect strategy for $\FMS$, respectively $\Gthree$, on faithful support contains four partial isometries $\beta_i$ from $f$ onto the four magic corners, with modular eigenvalue $r=2$, respectively $r=3$. In particular
\[
 \phi(m)=4r\phi(f),\qquad \phi(f)=\frac1{1+4r}.
\]
\end{lemma}
\begin{proof}
Choose a minimal $q$ of the centralizer $M_4$ and centralizer maps $u_x,v_y$ from $q$ to the calibrated rays. Their products are $u_x^*v_y=(-1)^{xy}q/\sqrt r$. On the standard corner Hilbert space $eJeJ L^2(M)$, where $e=f+q$, pull back the sign differences through $f+u_x$ and $J(f+v_y)J$. They are commuting self-adjoint contractions $X_x,Y_y$. The joint-support vector is
\[
 \theta_{xy}=(f+u_x^*)J(f+v_y^*)J\xi
             =(f+(-1)^{xy}q/\sqrt r)\xi.
\]
The rejected sign products give
\[
 X_x\theta_{xy}=(-1)^{xy}Y_y\theta_{xy},\qquad
 X_xY_y\theta_{xy}=(-1)^{xy}\theta_{xy}.
\]
Each coefficient $f\pm q/\sqrt r$ is invertible in the centralizer of the corner. Its vector separates both corner algebras; cancellation first makes all four contractions reflections. Put $H=f+q/\sqrt r$ and $Z=f-q$. The three positive cells identify one $X$ and one $Y$ with $XH\xi=YH\xi$. In the negative cell, $XZH\xi=-YZH\xi=-ZXH\xi$, since $Y$ commutes with $Z$. Separation gives $XZ+ZX=0$.

It follows that $X=\beta+\beta^*$, with $\beta=qXf$, $\beta^*\beta=f$ and $\beta\beta^*=q$. The self-adjoint mirror makes $X$ centralize $\omega(x)=\phi(kx)$, $k=f+q/r$. The bounded invertible centralizer perturbation formula of Pedersen--Takesaki \cite[Lemma 4.5]{PT} is
\[
 \sigma_t^\omega(x)=k^{it}\sigma_t^\phi(x)k^{-it}.
\]
Its $q$--$f$ block gives $\sigma_t^\phi(\beta)=r^{it}\beta$. Multiplication by the magic matrix units gives all four maps. KMS gives $\phi(q)=r\phi(f)$. If one sector vanished, the off-diagonal reflection would kill the other, contradicting $f+m=1$. This proves nonvanishing and the mass formula.
\end{proof}
This is the arbitrary-normal extension of the flagged-overlap mechanism in \cite[Appendix C]{FMS}. It is not obtained by substituting type-II multiplicities into a finite-dimensional self-test.

\subsection{Semifinite exclusion}
For a perfect $\GIII$ strategy, the first links make the flags $F_j$ a partition of the identity, with Bob mirrors $JF_jJ$. The second links identify the designated magic corner with $C_j=F_{j+1}+F_{j+2}$. Lemma~\ref{lem:bridge} supplies
\begin{equation}\label{eq:ck}
 w_j^*w_j=F_j,\quad w_jw_j^*=C_j,\quad
 \sigma_t^\phi(w_j)=2^{it}w_j,\quad
 \sum_jF_j=1,\quad\sum_jC_j=2.
\end{equation}
The mass equations $2\phi(F_j)=\phi(F_{j+1})+\phi(F_{j+2})$ imply $\phi(F_j)=1/9$ by a maximum principle around the cycle. Every copy has full support, since its flag and magic masses total $9\phi(F_j)=1$.

Suppose $M$ is semifinite and $\phi=\Tr(h\,\cdot)$ on faithful support. Then $hw_j=2w_jh$ as affiliated operators. With $v_j=w_j^*$, bounded functional calculus gives $g(h)v_j=v_jg(h/2)$. For a bounded nonnegative $g$ compactly supported in $(0,\infty)$, trace integration and the two covers in \eqref{eq:ck} yield
\[
 \phi(g(h))=\sum_j\Tr(hg(h)v_jv_j^*)
 =\frac12\sum_j\Tr(C_jhg(h/2))=\phi(g(h/2)).
\]
For rigor use $\Tr(T^{1/2}vv^*T^{1/2})=\Tr(v^*Tv)$ for positive integrable $T$, rather than cycling arbitrary unbounded products. The probability law of $h$ under $\phi$ would have equal mass on all dyadic windows $(2^k,2^{k+1}]$. Every such mass is zero, a contradiction.

For a semifinite Bob ambient algebra, standardize Bob first. His standard commutant is semifinite and contains the compressed Alice effects, so the same proof applies. Conditioning on a positive-mass semifinite central summand preserves zero losing events; hence all state weight of either ambient algebra lies on its type-III central part. An unused semifinite summand need not be absent. More generally, if Alice's algebra is contained in a semifinite separator $N\subseteq B'$ and the state is normal on $N$, perfect play is impossible.

\section{Prepared-state hosting and the flow of weights}\label{sec:hosting}
\subsection{Expected-inclusion transfer and explicit hosts}
\begin{lemma}[Transfer along expectations]\label{lem:transfer}
Let $E:M\to P$ be a faithful normal conditional expectation and let $\phi=\phi_P\circ E$ for a faithful normal state on $P$. Every standard-form correlation on $(P,\phi_P)$ is a correlation on $(M,\phi)$.
\end{lemma}
\begin{proof}
The space $K=\overline{P\xi_\phi}$ identifies with $L^2(P,\phi_P)$. Takesaki's expectation theorem \cite{Takesaki} makes $P$ modularly invariant. The orthogonal projection onto $K$ commutes with the modular objects, and the Tomita operators agree on $P\xi_\phi$. Thus $J_M|_K=J_P$. For $A,B\in P$,
\[
 \langle\xi_\phi,AJ_MB J_M\xi_\phi\rangle
 =\langle\xi_{\phi_P},AJ_PB J_P\xi_{\phi_P}\rangle.
\]
Embedding the left effects and using $J_MB J_M$ for the right effects preserves positivity, sums and every probability. The same identity transfers normal measurable POVMs for infinite games when those games and their integrals are defined; the finite-game compactness theorem below has a separate finiteness premise. Support compression handles nonfaithful starting marginals. Expectations compose.
\end{proof}
In the uniform KMS representation of the Cuntz algebra $\mathcal O_2$, let $s_i^*s_k=\delta_{ik}$, $\sum_i s_is_i^*=1$, and $\sigma_t(s_i)=2^{-it}s_i$. Its weak closure is $R_{1/2}$; these are established KMS and factor facts \cite{CK,Harris}. On $M_9\bar\otimes R_{1/2}$ take $\phi=\mathrm{tr}_9\otimes\varphi$. For each of the four pairs
\[
 \{j+1,j+2\},\ \{j+3,j+4\},\ \{j+5,j+6\},\ \{j+7,j+8\}
\]
put
\[
 F_a=e_{aa}\otimes1,\qquad W_{j;a,b}=e_{aj}\otimes s_1^*+e_{bj}\otimes s_2^*.
\]
Then $W^*W=F_j$, $WW^*=F_a+F_b$ and $\sigma_t(W)=2^{it}W$. With $U_0=F_j$ and $U_1,\ldots,U_4$ these four maps, $E_{ab}=U_aU_b^*$ are unital $M_5$ matrix units with weights $d=(1,2,2,2,2)$. Embed each real ideal effect $P$ as $A_P=\sum P_{ab}E_{ab}$ and Bob's as $JA_QJ$. Standard form gives
\begin{equation}\label{eq:block}
 \langle\xi,A_PJA_QJ\xi\rangle
 =\phi(F_j)\sum_{a,b}P_{ab}Q_{ba}\sqrt{d_ad_b}.
\end{equation}
Since $\phi(F_j)=1/9$, these are precisely the ideal probabilities of $\Psi_2$. The label measurement is $\{JF_aJ\}$. Both link partitions are centralizer partitions, so they match perfectly. Finite matrix amplification preserves the factor type, giving the claimed $R_{1/2}$ strategy.

There is an analogous explicit $R_{1/3}$ strategy for $\Gthree$. In $M_{13}\bar\otimes\pi_\varphi(\mathcal O_3)''$ with uniform KMS state set
\[
 F=e_{00}\otimes1,\qquad W_k=\sum_{i=1}^3e_{3(k-1)+i,0}\otimes s_i^*,\quad 1\le k\le4.
\]
The source is $F$ and the four orthogonal ranges partition $1-F$. Their eigenvalue is $3$, and the associated $M_5$ units have weights $(1,3,3,3,3)/13$. Formula~\eqref{eq:block} with denominator 13 gives the ideal $\Psi_3$ probabilities. The host identification is the uniform three-isometry case of the known Cuntz/KMS classification \cite[Theorem 4.9 and Example 4.13]{Harris}.

\subsection{The forced expected subfactor}
In any perfect $\GIII$ strategy let $V_j=w_j^*$ in \eqref{eq:ck}. They are a represented Cuntz--Krieger family for
\[
 A_{ja}=1\quad\Longleftrightarrow\quad a=j+1\text{ or }j+2.
\]
This nonpermutation graph is irreducible, every cycle has an exit, and every vertex projection is nonzero. The Cuntz--Krieger uniqueness theorem therefore makes the representation faithful \cite{CK}. The row sums are two and $\min(A^8)=1$. Return cycles of lengths five and nine give period one.

The modular action on this algebra is the gauge action with parameter $-t\ln2$. Its restricted state is the unique critical KMS state at inverse temperature $\ln2$, by \cite[Theorem 4.3]{Graph}. Let $P=W^*(V_j)$. Periodic Fourier averaging identifies its centralizer with the weak closure of the fixed AF algebra. Primitivity makes its normalized trace unique and extremal by Perron--Frobenius convergence. Matrix sizes grow without bound, so this centralizer is a diffuse finite factor. Any central element of $P$ is in that centralizer; hence $P$ is a factor.

The modular spectrum is $\{0\}\cup2^\Z$: admissible words supply every degree and periodicity excludes other positive values. Connes's factorial-centralizer criterion \cite[Corollary 3.2.7(a)]{Connes} identifies this with $S(P)$. Nuclearity gives injectivity of the weak closure, and its cyclic representation is separable. Injective-factor classification identifies $P$ with $R_{1/2}$ \cite{Connes76}. Takesaki gives the faithful normal state-preserving expectation $M\to P$.

This argument takes place on faithful support. In a factor with a nonfaithful hosting state, exclusion first makes the factor type III, and a nonzero type-III corner is isomorphic to the whole factor. Transporting the expected inclusion through such an isomorphism gives an expected $R_{1/2}$ in $M$ itself. Conversely Lemma~\ref{lem:transfer} transfers the explicit strategy along any such expectation. This proves the first two equivalences of Theorem~\ref{thm:host}.

\subsection{Halving on the core and the main equivalence}
In the continuous core, let $\lambda(t)=h^{it}$ and use its canonical trace $\tau$. For bounded nonnegative $g$ compactly supported in $(0,\infty)$ and central $z\ge0$, \eqref{eq:ck} gives the \emph{halving law}
\begin{equation}\label{eq:halving}
 \mu_g(z):=\tau(zg(h))=2\tau(zg(h/2)).
\end{equation}
The same positive-integrable trace identity used in the exclusion proves this without unbounded cyclic permutations. Its integrability is explicit:
\[
 \tau(g(h))=\phi(1)\int_0^\infty g(r)\,\frac{dr}{r^2}<\infty.
\]
At $s=\ln2$, trace scaling gives
\[
 \tau(\theta_s(z)g(h/2))=\tfrac12\tau(zg(h)).
\]
Applying \eqref{eq:halving} to $\theta_s(z)$ shows $\mu_g\circ\theta_{\ln2}=\mu_g$. Choose $g$ with positive integral and average over a period:
\[
 \nu=\frac1{\mu_g(1)\ln2}\int_0^{\ln2}\mu_g\circ\theta_s\,ds.
\]
This is a normal state invariant under the entire flow. Factor corners have isomorphic flows, so the conclusion passes from faithful support to the original factor.

We state the decisive analytic import exactly in the form needed: for a type-III factor with separable predual, an invariant normal state on its flow is equivalent to a normal expected embedding of $R_\infty$, and equivalent to such an embedding of every $R_\lambda$, $0<\lambda<1$ \cite[Theorem 3.5]{HM}. Lemma 3.4 there makes the expectation faithful after support reduction and type-III corner transport. No injectivity or bicentralizer assumption is required on the ambient factor. Combining this theorem, the halving law and transfer proves Theorem~\ref{thm:host}. Semifinite factors are excluded separately; their translation flow has no invariant normal probability state.

\subsection{Examples and the exact scope of blindness}
All separable-predual $\mathrm{III}_\lambda$ factors with $0<\lambda\le1$ host $\GIII$. Hosting does not imply $1/2\in S(M)$. Two examples also show directly how expected inclusions change subtype.

Put $Q=R_{1/2}$, $T=2\pi/\ln2$, and $\gamma=\sigma_{T/2}$. It is an outer order-two automorphism, since $\gamma(s_i)=-s_i$. The crossed product $L=Q\rtimes_\gamma\Z_2$ is a factor with its Fourier expectation onto $Q$. If $u$ implements $\gamma$ and $p=(1+u)/2$, then $pLp\cong Q^\gamma$. Every even monomial is generated by $s_is_j$: when both word lengths are odd, insert $\sum_k s_ks_k^*=1$. Thus the fixed algebra is the weak closure of the uniform $\mathcal O_4$ representation, hence $R_{1/4}$. The nonzero corner and $L$ are isomorphic. Therefore $R_{1/4}$ contains expected $R_{1/2}$ and hosts the game.

For a type-$\mathrm{III}_0$ example, let $t_0=2\pi/\ln3$ and $L_0=Q\rtimes_{\sigma_{t_0}}\Z$. Every nonzero power is outer, since $\ln2/\ln3$ is irrational, so Fourier coefficients prove factoriality. Amenability preserves injectivity. For the dual state $\psi$, $\sigma_T^\psi=\id$ and $\sigma_{t_0}^\psi=\operatorname{Ad}u$, where $u$ is in its centralizer. Choose bounded self-adjoint $A$ with $e^{iA}=u^*$ and put $k=e^{A/t_0}$. The perturbation by $k$ has period $t_0$. The two periodic states force $S(L_0)\subseteq\{0,1\}$. Hosting through the Fourier expectation excludes semifiniteness; hence $L_0$ is $\mathrm{III}_0$.

\begin{proposition}[Injective blindness]\label{prop:blind}
If $P$ is any injective factor with separable predual and $M$ is a factor in the hosting class of Theorem~\ref{thm:host}, every finite game hosted by $P$ is hosted by $M$. If $M$ is also injective, its hosted games are exactly those of $R_\infty$.
\end{proposition}
\begin{proof}
After support reduction, $P\subset P\bar\otimes R_\infty$ is expected by $\id\otimes\rho$. Haagerup--Musat's absorption lemma \cite[Lemma 3.8]{HM} identifies that tensor product with $R_\infty$. Theorem~\ref{thm:host} and Lemma~\ref{lem:transfer} finish the transfer. Apply it in both directions for the second assertion.
\end{proof}
In particular no prepared-state hosting test in this injective flow-finite class reads a nontrivial ratio-membership condition. This is not a statement about every noninjective host. The consequence of $\mathrm{MIP}^*=\mathrm{RE}$ recalled in \cite{JNVWY,Tensor} gives a finite synchronous game with a perfect tracial commuting model but spatial value at most $1/2$. Central decomposition of that model preserves its finitely many zero equations on almost every factor fibre. Choose such a separable-predual finite factor $P$. It cannot be injective, since injective finite-game correlations are in $C_{qa}$. The factor $P\bar\otimes R_\infty$ is $\mathrm{III}_1$ and hosts the game through the product expectation, while $R_\infty$ cannot. Thus noninjective flow-finite hosts do separate. Conversely a positive perfect-game witness cannot force injectivity: tensor any injective host with $L(\mathbb F_2)$ and use the product expectation.

\subsection{Finite shadows of the particular perfect correlation}
Truncate binary words at length $L$, and let $t_i$ prepend a letter, killing length-$L$ words. Use the diagonal state
\[
 h_{a,\mu}=\frac{2^{-|\mu|}}{9(L+1)}.
\]
Replacing $s_i^*$ by $t_i^*$ gives exact five-dimensional blocks on the good subspace. In each copy the complement is the flag empty-word atom and the eight magic top-level atoms; its state mass is $1/(L+1)$. Complete these atoms by deterministic strategies losing just one internal cell and preserving both links. Explicitly the flag strategy always flags on rows/columns and answers $+$ on signs. For magic completions use the following row/column triples, with $T=(+,C_1)$ and $U=(+,R_1)$:
\[
\begin{array}{c|ccc|ccc}
 &R_1&R_2&R_3&C_1&C_2&C_3\\\hline
\text{designated}&+++&+++&+++&+++&+++&++-\\
\text{other}&+++&+--&+--&+++&+--&+-+
\end{array}
\]
They lose only $(R_3,C_3)$; the designated case is used precisely on atoms $j+1,j+2$. Direct substitution verifies every rule. It follows that
\[
 \ell_L=\frac1{27(L+1)},\quad
 S_L=\ln9+\ln(L+1)+\tfrac L2\ln2,\quad
 d_L=9(2^{L+1}-1).
\]
Convergence of scores alone is insufficient for a claim about a particular correlation. Each ideal effect is a fixed degree-one polynomial in the Cuntz letters, with $9\times9$ coefficients. Boundary corrections have vanishing state mass. The truncated moments are
\[
 \phi_L(t_i^*t_j)=\delta_{ij}\frac L{L+1},\qquad
 \phi_L(t_it_j^*)=\delta_{ij}\frac L{2(L+1)},\qquad\phi_L(t_i)=0.
\]
Together with exact modular scaling on these shifts they give convergence of every left/right polynomial correlation. Cauchy--Schwarz removes the boundary corrections uniformly for this finite list of effects. Hence all question pairs, including weight-zero cross-copy pairs, converge to the displayed KMS correlation. This proves its membership in $C_{qa}$. Exclusion puts it outside $C_{qs}$ and proves nonattainment, while all relevant value suprema are one.

\section{Fixed-state gaps and the ratio invariant}\label{sec:states}
\begin{lemma}[Finite-game spectral obstruction]\label{lem:gap}
Suppose a finite game forces a nonzero modular $r$-eigenoperator in every faithful normal perfect realization on every $\sigma$-finite standard-form algebra. If $r\notin\Sp\Delta_\phi$, then $v_G(M,\phi)<1$.
\end{lemma}
\begin{proof}
Assume fixed-state losses tend to zero. Put all coordinate effects on $L^2(M,\phi)$ and form the Hilbert ultraproduct with constant vector $\Omega=(\xi_\phi)_\omega$. The Groh--Raynaud algebra $N$ is standard and its commutant is the ultraproduct of the commutants \cite[Theorems 3.18 and 3.22]{AH}. Coordinate effects are POVMs, and every positively weighted losing event has zero ultralimit, because its probability is bounded by total loss divided by its weight.

Let $p$ be the support of $\Omega$ in $N$, $q=J_\omega pJ_\omega$, and $e=pq$. Compress \emph{both} parties to $eH_\omega$. Since $p$ commutes with Bob and $q$ with Alice, these commuting POVMs preserve all correlations. The standard compressed left algebra is the Ocneanu ultrapower and the vector state is faithful \cite[Theorem 3.7 and Corollary 3.28]{AH}. These statements do not assume arbitrary effect sequences normalize the Ocneanu ideal.

Forcing in this possibly nonfactor algebra yields $x\ne0$ with modular eigenvalue $r$. Hence $x\xi_{\phi^\omega}\ne0$ is an $r$-eigenvector of its modular operator. But the constant-pair spectrum equality \cite[Corollary 4.8(3)]{AH} is
\[
 \Sp\Delta_{\phi^\omega}=\Sp\Delta_\phi.
\]
This contradicts the premise. The equality concerns spectra, not point spectra. The coordinate modular action is identified by Theorem 4.1 of the same source.
\end{proof}
Finiteness is essential to this compactness statement. For example, let Alice answer $n\ge1$, let Bob receive $k$ with probability $2^{-k}$ and have one answer, and accept iff $n\ge k$. Value one is approached by $n\to\infty$, while no perfect normal countably additive strategy exists even on scalars. Thus a vacuous perfect-play forcing premise for an infinite-answer game gives no spectral obstruction.

Connes's state-intersection result \cite[Corollary 3.2.5(d)]{Connes} says that failure of a positive ratio in $S(M)$ is witnessed by some faithful normal state's modular spectrum. The possible extra zero for properly infinite semifinite factors is irrelevant. Lemmas~\ref{lem:bridge} and \ref{lem:gap} therefore prove necessity of $1/2\in S(M)$ for $\FMS$ and of $1/3\in S(M)$ for $\Gthree$. The same necessity holds for $\GIII$.

\subsection{Sufficiency and a phase-compensated state approximation}
We record the continuity estimate used throughout. If effects have zero loss in a normal state $\psi$, their positive loss contraction $L$ annihilates $\xi_\psi$. Keeping the effects gives
\begin{equation}\label{eq:continuity}
 \langle\xi_\phi,L\xi_\phi\rangle
 \le\|\xi_\phi-\xi_\psi\|^2\le\|\phi-\psi\|.
\end{equation}
The last inequality is the natural-cone inequality. It also applies after the standard unitary implementer $uJuJ$ conjugates both measurement families. Nonfaithful intermediate states cause no problem: extend their support strategies by arbitrary POVMs outside the support.

In type $\mathrm{III}_1$, a faithful product state on $M_5\bar\otimes N\cong M$ with marginal $\rho_r=\operatorname{diag}(1,r,r,r,r)/(1+4r)$ hosts the ideal $r=2$ or $r=3$ game. Connes--St\o rmer approximate inner homogeneity \cite{CS,Haagerup} makes its inner orbit norm dense among normal states. Formula~\eqref{eq:continuity} proves value one at every faithful normal state. For $\GIII$, Theorem~\ref{thm:host} gives a faithful hosting state by composing the expected $R_{1/2}$ state with its faithful expectation, and the same argument works.

For completeness, the $\mathrm{III}_\lambda$ direction is proved using one state-orbit argument for both ratios. This is the phase-compensated spectral-state construction; a direct dual-density rounding construction gives an alternative. We use the following named imports from Haagerup--St\o rmer \cite{HS}. The spectral state
\[
 \widehat\eta(z)=\tau\bigl(z1_{(1,\infty)}(h_\eta)\bigr)
\]
on the center of one fixed continuous core classifies approximate unitary equivalence, with
\begin{equation}\label{eq:orbit}
 \inf_{u\in\mathcal U(M)}\|\phi-u\psi u^*\|=\|\widehat\phi-\widehat\psi\|.
\end{equation}
For $\mathrm{III}_\lambda$, put $L=-\ln\lambda$. Proposition 10.5 and its proof identify permitted spectral states with mixtures of translates on $\R/L\Z$ of the normalized kernel
\[
 K(y)=e^{-y}/(1-e^{-L}),\qquad0\le y<L,
\]
with norm-continuous translation. This use of a circle orientation fixes all signs below.

The discrete decomposition of an existing $\mathrm{III}_\lambda$ factor is the direction of \cite[Theorem 4.4.1(b)]{Connes}. Its generalized trace $\omega$ has period $2\pi/L$ and factorial centralizer $N$ of type $\mathrm{II}_\infty$. If $e\in N$ has finite positive trace $t=\omega(e)$, the state $\psi_e(x)=\omega(exe)/t$ has a factorial $\mathrm{II}_1$ centralizer on $eMe$. Suppose $r=\lambda^{-k}$ with integer $k\ge1$. The $r$-homogeneous space is a nonzero full bimodule over the centralizer. For source and target centralizer projections of masses $a$ and $ra$, take a maximal orthogonal family of homogeneous partial isometries. KMS gives the same mass ratio for each member. If a remainder exists both remainders are nonzero; factorial comparison and fullness produce another homogeneous map between subprojections, a contradiction. The bounded strong sum has exactly the specified source and range.

For $r=3$ choose a partition with masses $(1,3,3,3,3)/13$ and four such maps. Formula~\eqref{eq:block} supplies $\Gthree$. For $r=2$ the same construction supplies $\FMS$; choosing nine equal atoms and four disjoint target pairs for each supplies $\GIII$. Thus \emph{every} finite generalized-trace reduction hosts the required game.

Two spectral-state identities control the approximation:
\begin{equation}\label{eq:scaling}
 \widehat{a\eta}=a\widehat\eta\circ\theta_{\ln a},\qquad
 \widehat{\sum_i\eta_i}=\sum_i\widehat\eta_i\quad\text{for orthogonal supports}.
\end{equation}
The first follows from $h_{a\eta}=ah_\eta$, the threshold projection, and trace scaling. The second follows from orthogonality of the affiliated supports. No arbitrary affinity is asserted. Finite reductions have
\[
 \widehat{\psi_e}=K\circ\theta_{-\ln t}.
\]
Indeed $e\mapsto\tau(ze1_{(1,\infty)}(h_\omega))$ is a tracial normal weight on the factorial $N$, hence proportional to $\omega|_N$. Normalize at trace one and apply \eqref{eq:scaling}. The shape of $K$ follows from the geometric spectrum $\lambda^\Z$ in the compact modular crossed product and its trace-scaling shift; on the fundamental interval the distribution is constant before passing to logarithmic coordinates. This is also the periodic-state kernel in the proof of the cited proposition.

Approximate $\widehat\phi$ in norm by a finite mixture $\sum_i a_iK\circ\theta_{b_i}$. Choose orthogonal finite projections $e_i\in N$ with traces $t_i>0$ satisfying
\[
 -\ln t_i+\ln a_i=b_i\pmod L.
\]
Any finite list of positive traces is available in $\mathrm{II}_\infty$. The block state $\psi=\sum_i a_i\psi_{e_i}$ hosts the game, and \eqref{eq:scaling} gives exactly the desired spectral mixture. The coefficient-dependent phase shift has thus been compensated. Equation~\eqref{eq:orbit}, followed by \eqref{eq:continuity}, approaches $\phi$ with hosting states and proves $v_G(M,\phi)=1$.

We have proved all the asserted worst-state criteria, including
\[
 \vworst_{\GIII}(M)=1\iff 1/2\in S(M).
\]
For $r=2$ the positive sufficient subtypes are $\mathrm{III}_{2^{-1/k}}$ and $\mathrm{III}_1$; for $r=3$ they are $\mathrm{III}_{3^{-1/k}}$ and $\mathrm{III}_1$. No injectivity hypothesis enters this proof.

\subsection{Two games, one-state obstruction, and algebraicity}
If both worst-state values are one, $2$ and $3$ belong to the closed group $S(M)\setminus\{0\}$. Unique factorization makes $\ln2/\ln3$ irrational, so $\Z\ln2+\Z\ln3$ is dense. The group is all $\R_{>0}$, and $M$ is $\mathrm{III}_1$. The converse was proved by homogeneity. In fact the two-game conclusion needs only the fixed-state obstruction and the $\mathrm{III}_1$ argument, independently of the more involved $\mathrm{III}_\lambda$ approximation.

Nevertheless the faithful density $\rho_2\otimes\rho_3$ on $M_{25}$ wins both games after purification, by playing each on its own tensor factor. Its eigenvalues are $1,2,3,6$ divided by 117, with multiplicities $1,4,4,16$. Here $S(M_{25})=\{1\}$. The two-game property quantifies separately over game, state and accuracy; neither common measurements nor a common flag are required. A finite set of preparations cannot establish that premise.

\begin{proposition}[Limited algebraicity]\label{prop:alg}
If a finite binary-scoring game has a finite-dimensional perfect strategy and universally forces a fixed modular ratio $r$, then $r$ is algebraic.
\end{proposition}
\begin{proof}
Fix dimensions in which perfection occurs and purify after a finite enlargement. Real and imaginary coordinates for the vector and POVMs, positive factorizations $E=C^*C$, normalization, completeness and zero losing probabilities form a finite existential system over $\Q$. Only the support of the question distribution matters. The real algebraic numbers form a real closed subfield, so real-closed-field transfer gives an algebraic solution. The forced ratio is not inserted as a coefficient of this system. The reduced density has algebraic entries and algebraic nonzero eigenvalues. On faithful support its type-I modular eigenvalues are ratios of these eigenvalues. Universal forcing applies to this algebraic realization, so $r$ is algebraic.
\end{proof}
This does not settle algebraicity when all perfect realizations require infinite dimension. Harris's minimal exact-embezzling-factor algebraicity theorem has different hypotheses \cite{Harris}. A finite noncommutative presentation alone supplies no finite-dimensional real-closed-field argument.

\section{Which ratios finite games can force}\label{sec:ratios}
\begin{definition}[Forced ratio]\label{def:forced}
A finite binary-scoring game \emph{forces} $r>0$ if it has at least one perfect normal strategy and every perfect normal strategy, after the support reduction of Lemma~\ref{lem:support}, has a nonzero bounded $r$-eigenoperator for its faithful marginal state. The ambient algebra need not be a factor. The identity and adjoints supply ratio $1$ and reciprocals. A distinguished bridge parameter is not the entire forced set: the ideal finite-dimensional $\FMS$ strategy has modular point spectrum $\{1,2,1/2\}$.
\end{definition}

\subsection{The general overlap bridge}
A rigid real core of dimension $d$ means checked base measurements consisting of a common flag $f$ and rank-one ray projections in a centralizer copy of $M_d$ with unit $m=1-f$, with arbitrary multiplicity. On every base/base pair accept two flags, reject exactly one flag, and accept two ray labels precisely when their ordinary dot product is nonzero. The six magic-square contexts give such a core for $d=4$, by Section~\ref{sec:games}.

Choose four forced unit rays, with phases such that
\[
 \langle w_x,v_y\rangle=c(-1)^{xy},\qquad x,y\in\{0,1\},\quad c>0.
\]
The rays within either pair may be nonorthogonal and may occur in different base questions. Add $T_x$ for Alice and $U_y$ for Bob. Each bridge question replaces the flag and its selected ray in a calibrating basis by two sign answers, retaining the other $d-1$ ray labels. Calibrate $T_x$ against Bob's basis of $w_x$: a sign matches exactly the flag or selected ray, and a residual answer matches its own label. Use the reversed rule for $U_y$ and Alice's basis of $v_y$. At $(T_x,U_y)$ two signs must satisfy $ab=(-1)^{xy}$; otherwise accept. All other cells are free. Give every question pair positive weight.

\begin{lemma}[Arbitrary-normal flagged overlap]\label{lem:overlap}
This bridge forces $r=c^{-2}\ge2$. On faithful support it extracts a partial isometry $\beta\in qMf$, where $q$ is a minimal diagonal of the rigid core, satisfying
\[
 \beta^*\beta=f,\qquad \beta\beta^*=q,\qquad
 \sigma_t^\phi(\beta)=r^{it}\beta,\qquad
 \phi(f)=\frac1{1+dr}.
\]
Both supports are nonzero; ``full'' refers to these initial and final supports.
\end{lemma}
\begin{proof}
Choose centralizer maps $u_x=|w_x\rangle\langle e|$ and $v_y=|v_y\rangle\langle e|$, so $u_x^*v_y=c(-1)^{xy}q$. Put $e_0=f+q$ and work on $e_0Je_0J L^2(M)$ with separating vector $e_0\xi$. Calibration fixes the sign supports. Pulling their self-adjoint contraction differences back by $f+u_x$ and $J(f+v_y)J$ gives commuting $X_x,Y_y$ and
\[
 \theta_{xy}=(f+u_x^*)J(f+v_y^*)J\xi
       =(f+c(-1)^{xy}q)\xi.
\]
Forbidden signs imply $X_x\theta_{xy}=(-1)^{xy}Y_y\theta_{xy}$ and $X_xY_y\theta_{xy}=(-1)^{xy}\theta_{xy}$. Positivity of contraction defects makes their squares fix $\theta_{xy}$. Each $f\pm cq$ is invertible in the corner centralizer. Its vector separates the left algebra by cancellation and the right algebra by commuting this coefficient through the commutant. Thus all four contractions are reflections.

Put $H=f+cq$ and $Z=f-q$. The three positive cells identify reflections $X,Y$ with $XH\xi=YH\xi$; the negative cell gives $(XZ+ZX)H\xi=0$. Separation yields $XZ+ZX=0$. Consequently $X=\beta+\beta^*$ with $\beta=qXf$ and the asserted supports. The mirror relation makes $X$ centralize the faithful corner functional $\omega(x)=\phi(kx)$, $k=f+c^2q$. The bounded centralizer perturbation formula \cite[Lemma 4.5]{PT} gives
\[
 \sigma_t^\omega(x)=k^{it}\sigma_t^\phi(x)k^{-it},\qquad
 \sigma_t^\phi(\beta)=k^{-it}\beta k^{it}=c^{-2it}\beta.
\]
KMS gives $\phi(q)=c^{-2}\phi(f)$. The $d$ core diagonals are centralizer equivalent, proving the mass formula. If one support vanished the off-diagonal reflection would kill the other and then $m+f=1$ would vanish. Finally Cauchy--Schwarz gives $4c^2\le\|w_0+w_1\|^2$ and $4c^2\le\|w_0-w_1\|^2$; adding proves $c^2\le1/2$.

Completeness uses
\[
 \psi_c=\frac{c|ff\rangle+\sum_{i=1}^d|ii\rangle}{\sqrt{c^2+d}},
 \quad t_x^a=\frac{f+aw_x}{\sqrt2},\quad
 u_y^b=\frac{f+bv_y}{\sqrt2}.
\]
The two-sign amplitude is $c(1+ab(-1)^{xy})/(2\sqrt{c^2+d})$; all other forbidden amplitudes are calibration or base orthogonalities. No within-pair orthogonality was used.
\end{proof}
The overlap mechanism, its nonorthogonal permission and ideal state are due to Man\v cinska \cite[Appendix C]{FMS}. Her Theorem~17 gives the finite-dimensional state self-test for a suitable base game; the lemma supplies the arbitrary-normal extension.

\begin{proposition}[Pauli confinement]\label{prop:pauli}
For stabilizer projections $P_A,P_B$, every nonzero eigenvalue of $P_AP_BP_A$ is $2^{-k}$ for an integer $k\ge0$. The flagged constructions generated homogeneously by such bridge tilts force only ratios in $2^\Z$.
\end{proposition}
\begin{proof}
For stabilizer groups $S_A,S_B$, put $H=S_B\cap C(S_A)$. Terms anticommuting with some element of $S_A$ vanish under compression. Hence
\[
 P_AP_BP_A=\frac{|H|}{|S_B|}P_AP_H,
 \qquad P_H=\frac1{|H|}\sum_{g\in H}g.
\]
$P_H$ is a projection commuting with $P_A$, so their product is a projection, possibly zero. Both group orders are powers of two. No single-character action on a higher-rank code is assumed. The last assertion follows from Proposition~\ref{prop:degree} below. Non-stabilizer rays on the same qubit core escape this restriction.
\end{proof}

\subsection{All rational rays on the binary magic-square core}
A rational ray has a nonzero rational representative; its unit representative need not be rational. The six seed contexts contain 24 rays: four coordinate rays, twelve pair sums and differences, and eight sign rays. Use the following basis rule: add a rational orthogonal basis if, for each new ray $b$, previously forced rays perpendicular to $b$ span $b^\perp$.

The rule is sound in arbitrary normal strategies. Checked matching makes the new answers centralizer projections with the same flag. For any old perpendicular ray $u$, the forbidden positive expectation and faithfulness imply $E_bP_u=0$. Spanning gives $E_b\le P_b\otimes1$. Completeness makes the four positive differences sum to zero, so equality holds for each answer. This proof also works directly with initially unsharp effects after the checked partition reduction.

\begin{theorem}[Rational-ray closure]\label{thm:closure}
The least closure of the 24 seed rays under the basis rule is every rational ray of $\R^4$. Every rational orthogonal basis can therefore be forced by a finite chain of added questions.
\end{theorem}
\begin{proof}
Identify $\R^4$ with the quaternions. Signed coordinate permutations and left multiplication by $(1+i)/\sqrt2$, $(1+j)/\sqrt2$, $(1+k)/\sqrt2$ permute the seed rays and preserve rational rays. Applying a symmetry to every question of an existing finite chain proves invariance of the closure; this is an actual image chain.

Induct on the squared norm $n$ of a primitive integral representative $w$. For even $n$, primitiveness gives two or four odd coordinates. Some $a\in\{i,j,k\}$ pairs equal-parity positions, so $u=(1-a)w/2$ is integral with norm squared $n/2$, and $w=(1+a)u$. Induction and seed symmetry apply; primitive reduction only lowers the norm.

For a vector with at least three nonzero coordinates, fix $w_i\ne0$. The three vectors $w_je_i-w_ie_j$, $j\ne i$, are independent in $w^\perp$ and have squared norms strictly below $n$. Apply induction also to the signed permutations $wi,wj,wk$. The quaternion orthogonal basis $\{w,wi,wj,wk\}$ now meets the basis rule.

For a planar vector $(p,q,0,0)$ with $pq\ne0$, the even case is done. If $n$ is odd, the vectors $(q,-p,1,0)$ and $(p,q,1,0)$ have even norm squared $n+1$ and descend to $(n+1)/2<n$. Together with $e_2,e_3$ they span the respective complements of $(p,q,0,0)$ and $(q,-p,0,0)$. Thus their planar orthogonal basis is forced. One-coordinate rays and norms one and two are seeds. Finally combine finitely many ray chains to force any particular rational orthogonal basis. This proves finite existence without a game-size bound.
\end{proof}

\subsection{Single bridges, full-support products and positive rationals}
\begin{theorem}[Orthogonal-pair bridge parameters]\label{thm:orthogonal}
The distinguished parameters of rational rank-one bridges in $\R^4$ with both within-pair orthogonalities are exactly
\[
 \mathcal R_4=\{2+x^2+y^2:x,y\in\Q\}.
\]
Every member is forced on the binary magic-square core.
\end{theorem}
\begin{proof}
Equal nonzero normalized overlaps imply that the squared norms of rational representatives $w_0,w_1$ have rational-square ratio. Rescale them to common squared norm $a$ and orient them to have equal dot products against $v_0$. A rational quaternion frame $\{w_0,w_0i,w_0j,w_0k\}$ has Gram matrix $aI$. If $w_0i\ne w_1$, the rational Householder reflection with normal $w_0i-w_1$ fixes $w_0$ and sends $w_0i$ to $w_1$; otherwise use the identity. In that frame write $v_0=(t,t,z,h)$ with rational coefficients and $t\ne0$. Then $c^2=t^2/(2t^2+z^2+h^2)$, proving necessity even with only the $w$-pair orthogonal.

Conversely, for rational $\alpha,\gamma$, choose $w_0=e_0,w_1=e_1$ and representatives
\[
 v_0=(1,1,\alpha,\gamma),\qquad v_1=(1,-1,-\gamma,\alpha).
\]
Both pairs are orthogonal and their normalized cross overlaps are $(-1)^{xy}(2+\alpha^2+\gamma^2)^{-1/2}$ for bridge indices $x,y\in\{0,1\}$. Apply Theorem~\ref{thm:closure} and Lemma~\ref{lem:overlap}.
\end{proof}
The classical two-square criterion says that a nonzero positive rational is a sum of two rational squares exactly when its valuation at each prime congruent to $3$ modulo $4$ is even, including denominator valuations. Necessity follows by clearing denominators; such a prime dividing a sum of two integer squares divides both entries. Sufficiency uses the classical prime two-square theorem \cite{Zagier}, multiplicativity and inversion. Thus $5,8,9,13,16,17$ are excluded as orthogonal-pair single-bridge parameters, with no exclusion asserted for products or nonorthogonal pairs.

The exact nonorthogonal computations in \texttt{anc/ratios/} supply witnesses for $15/4$, $91/16$, $20/3$, $15/2$, $27/2$, $14$, $30$, $90$, $63/4$ and $39/2$. These are finite computations only. For example take
\[
 w_0=(1,2,-3,0),\quad w_1=(1,2,0,3),\quad
 v_0=(1,0,0,0),\quad v_1=(1,-2,-2,0).
\]
The squared norms are $14,14,1,9$ and cross dot products $(1,3;1,-3)$, so $r=14$. Since $14-2=12$ fails the two-square criterion, this lies outside $\mathcal R_4$. No exact description of all nonorthogonal parameters is established; a prospective Gram-form criterion must include singular forms and rank-three configurations.

\begin{lemma}[Full-support product]\label{lem:product}
The parallel product of two rigid full-support bridge games with parameters $r,r'$ forces $rr'$ and $r/r'$. Its ideal finite-dimensional strategy exists. Finite iteration has the same property.
\end{lemma}
\begin{proof}
Fix the other component's questions and take component answer marginals. Every resulting strategy is perfect. Checked base classes have the same Bob mirrors for all other-component questions, so separation makes them independent of those questions. For bridge sign differences the common relation $XH\xi=YH\xi$ with invertible $H$ gives the same independence; calibration handles the residual outcomes. Do this for both components and both parties.

All these marginal effects are projections. A joint POVM with sharp marginals has commuting marginals and joint effects their products: each effect lies below both marginals, and the positive residual defects sum to zero. Thus every extracted operator of one component commutes with those of the other, even for initial joint POVMs.

The core matrix units and $\beta$ generate a unital $M_{d+1}$ with unit $f+m=1$. The other component similarly generates $M_{d'+1}$. In its tensor decomposition $f=1\otimes z$, $z\ne0$, so its product with any nonzero minimal diagonal of that matrix algebra is nonzero. In particular
\[
 (\beta\beta')^*(\beta\beta')=ff'\ne0,\qquad
 (\beta\beta'^*)^*(\beta\beta'^*)=fq'\ne0.
\]
Their modular degrees are $rr'$ and $r/r'$. Tensor the ideal strategies for completeness. Finite products still extract a unital full matrix algebra with nonzero minimal product diagonals. This is not a product theorem for arbitrary eigenoperators in arbitrary games.
\end{proof}

\begin{theorem}[Every positive rational is forced]\label{thm:rationals}
For each $r\in\Q_{>0}$ there is a finite binary-scoring game with a finite-dimensional perfect strategy which forces $r$ in every perfect normal strategy.
\end{theorem}
\begin{proof}
Ratio two is supplied by $\FMS$. For an odd prime $p$, the residue sets of $a^2$ and $-2-b^2$ each have $(p+1)/2$ elements, so intersect. Choose $|a|,|b|\le(p-1)/2$. Then
\[
 2+a^2+b^2=pk,\qquad 0<k<p,
\]
since $2+(p-1)^2/2<p^2$. The left side belongs to $\mathcal R_4$. By prime induction every prime divisor of $k$ has a full matrix bridge-product construction; Lemma~\ref{lem:product} divides these factors out and forces $p$. Finite products, powers, adjoints and the identity give every positive rational.
\end{proof}
The game can self-test a tensor product of bridge states. This is existence of a game forcing each rational, not a self-test of every prescribed rational Schmidt spectrum. More precisely, this section partly answers Man\v cinska's Open Question~2 \cite{FMS} by the family of bridge spectra from $\mathcal R_4$ and the checked nonorthogonal witnesses, together with tensor products. Prescribing all rational spectra remains a different problem. Her CK-31 example already has $c=2/5$ and squared ratio $25/4$; her counting obstruction already limits perfect-game self-test spectra to a countable set.

\subsection{Six explicit games and exact certificates}
The six games for $r=3,5/2,6,7,10,25/4$ are specified in \texttt{anc/ratios/} by integer ray bases, ordered forcing certificates, selected oriented bridge rays and the rules above. Each question has five answers and the distribution is uniform on every question pair. The ideal state is $(|ff\rangle+\sqrt r\sum_{i=0}^3|ii\rangle)/\sqrt{1+4r}$. The checker constructs every measurement and predicate entry and tests all forbidden amplitudes and every pair's normalization by exact multiquadratic arithmetic. It also checks each spanning triple against strictly earlier rays. These are explicit finite examples, with no minimal-size claim; the ratio-three example here uses a different added basis from the conference game $\Gthree$ of Section~\ref{sec:games}.

\begin{center}
\begin{tabular}{rrrrrr}
\toprule
$r$ & $c^2$ & Questions/party & Pairs & Events & Forbidden\\
\midrule
$3$ & $1/3$ & 9 & 81 & 2025 & 732\\
$5/2$ & $2/5$ & 13 & 169 & 4225 & 1638\\
$6$ & $1/6$ & 9 & 81 & 2025 & 732\\
$7$ & $1/7$ & 13 & 169 & 4225 & 1642\\
$10$ & $1/10$ & 13 & 169 & 4225 & 1638\\
$25/4$ & $4/25$ & 17 & 289 & 7225 & 2902\\
\bottomrule
\end{tabular}
\end{center}

For that example the added basis is
\[
 (0,1,-2,-1),\quad(0,1,1,-1),\quad
 (1,-1,0,-1),\quad(2,1,0,1).
\]
Use $w_0=e_3,w_1=-e_1$ and $v_0=(0,-1,-1,1)/\sqrt3$, $v_1=(-1,1,0,1)/\sqrt3$. They have the required sign pattern, giving nine questions per player and 732 forbidden events. For ratio six use the norm-six pair of the same basis, with phases chosen for that pattern. The data files specify all other bases and calibrations without numerical thresholds.

\subsection{Degree bounds and the scope of arithmetic restrictions}
\begin{proposition}[Homogeneous-generation degree bound]\label{prop:degree}
Suppose one actual perfect witness strategy has faithful state $\phi$, and $P=W^*(\text{Alice's effects})$ is generated by modular homogeneous elements with degrees in a subgroup $L\subset\R_{>0}$, including adjoints. Every ratio forced by the game lies in $L$, without taking its topological closure.
\end{proposition}
\begin{proof}
$P$ is modular invariant and has the $\phi$-preserving conditional expectation \cite{Takesaki}. Equivalently restrict to $[P\xi]$ and compress Bob's effects; this preserves correlations and gives the faithful restricted state with the restricted modular group. Finite words in the homogeneous generators give a dense span of modular eigenvectors with degrees in $L$. A vector of any other degree is orthogonal to that entire span. A nonzero extracted eigenoperator applied to the separating vector cannot have that degree. Only this one witness is required.
\end{proof}

\begin{corollary}[Finite rigid cores and rational seeds]\label{cor:seeds}
Consider a finite collection of finite-dimensional exactly rigid cores and full-support bridges, with added effects admitting the checked homogeneous decomposition below. Every forced ratio lies in the multiplicative group of the bridge parameters, and those parameters are algebraic. If the cores start from rational seeds and use rational orthogonal-complement closure, all forced ratios are rational. The union over this rational construction class is exactly $\Q_{>0}$.
\end{corollary}
\begin{proof}
Checked class projections and core matrix units have degree one. A bridge sign difference splits into $u_x\beta$ and its adjoint, of degrees $r_i,r_i^{-1}$; these terms are themselves compressions of Alice's observables. Checked label and boundary projections and finite products preserve the stated decomposition when they are all the added effects. Proposition~\ref{prop:degree} applies whenever an actual perfect witness has these generators.

In each fixed finite core dimension, the finite projection, orthogonality and completeness equations, with trace-one rank conditions, are a polynomial system over $\Q$. Real-closed-field transfer gives a real algebraic solution; exact rigidity fixes its squared overlaps to those of the original core. Thus $c_i^2$ is algebraic. This argument concerns the finite core, even if a later glued game has only infinite-dimensional perfect realizations. It does not apply finite-dimensional feasibility to that glued strategy. Rational linear systems preserve rational rays under orthogonal-complement operations, so rational seeds give rational $c_i^2$. The degree bound and Theorem~\ref{thm:rationals} give the last assertion.
\end{proof}
This includes the specified type-III-only checked gluing when its witness has precisely that homogeneous generation, as in \eqref{eq:ck}. It asserts nothing for arbitrary extra gluings or infinite-dimensional rigid cores.

\begin{proposition}[Totally-real single-bridge rationality]\label{prop:totallyreal}
For a single rigid bridge whose ray representatives lie in a totally real number field, the distinguished parameter forced by that bridge is rational. The special ideal modular spectrum $\{1,r,r^{-1}\}$ is essential.
\end{proposition}
\begin{proof}
Real algebraic PVM strategies satisfy polynomial equations over $\Q$: symmetry, idempotence, completeness, state normalization and zero forbidden probabilities. Any real field embedding preserves them, hence gives another perfect strategy. When transporting density eigenvalues, enlarge the coefficient field to contain them and use a compatible extension of the embedding: the image density is positive semidefinite and real symmetric, so such real extensions exist. Images of square roots cannot be prescribed independently.

For a single bridge over a totally real field $K$, squared norms and $c^2$ are totally positive. Under every embedding the four nonzero cross dot products have product equal to a negative square, so phases can restore the frustrated pattern. The overlap bound gives $\tau(c^2)\le1/2$. Adjoining the needed square roots of norms, $c^2$ and $c^2+d$ preserves total reality. Thus the ideal strategy transports under every embedding of $K$, and its modular spectrum becomes $\{1,\tau(r),\tau(r)^{-1}\}$ with $\tau(r)\ge2$. Universal forcing of the original $r\ge2$ requires $r=\tau(r)$ for every embedding. Therefore $r\in\Q$.
\end{proof}
Transport for a general strategy gives a constraint on an entire spectrum, not rationality of every forced ratio. An arbitrary real algebraic field need not be totally real, so this proposition is not an exclusion of irrational ratios for general finite games.

Universally forced ratios give necessary spectral obstructions to worst-state value one by Lemma~\ref{lem:gap}; the rational construction theorem supplies no converse for an arbitrary forcing game. No theorem here reduces all finite-game type distinctions to universal ratio forcing.

\section{The price of approximate play}\label{sec:prices}
All entropies in this section use natural logarithms. Each optimized cost at accuracy $\eps$ is an infimum over the specified strategies with loss at most $\eps$. For type I, $S(h)=-\Tr(h\ln h)$ is marginal entropy. For a normalized finite trace $\tau(1)=\tau(h)=1$, the deficit is $D_\tau(h)=\tau(h\ln h)$. For a chosen semifinite trace let $T$ have law $p(t)dt$, $p(t)=\mu_t(h)$, and put $X=\log_2T$. Trace rescaling translates $X$ and leaves its variance and $H$ unchanged in a factor. In a nonfactor, global rearrangement can lose central information; the inequality below still holds for the chosen trace.

\subsection{A robust gluing theorem}
Theorem 7.2 of the companion base-game paper \cite{Base} supplies a universal finite constant $C_F$ for the flagged sector estimate, stated in the form needed here. For a projective spatial realization of one FMS copy, with flag $F$ and unweighted internal forbidden mass $N$, its sorted Schmidt-amplitude lists satisfy
\begin{equation}\label{eq:sector}
 \bigl\|\operatorname{Sch}((1-F)\psi)
       -4\mathbin{\times}\sqrt2\operatorname{Sch}(F\psi)\bigr\|_2
 \le C_F\sqrt N.
\end{equation}
Here $k\mathbin{\times}\Lambda$ means the sorted union of $k$ copies of a list, padded by zeros. The estimate is homogeneous in unnormalized vectors, includes arbitrary multiplicity, and uses the unweighted sum of the 25 forbidden masses. The companion proof uses the order-32 Pauli-group extraction and bounded words of length at most six; its estimates use no finite-rank cutoff. We do not give a validated numerical upper bound for $C_F$ here. We give the additional nine-copy gluing argument in full.

Purify a mixed strategy on Bob's side; a common Naimark dilation with pure input registers makes the finitely many POVMs projective without changing Alice's nonzero spectrum. Let $\Lambda=\operatorname{Sch}(\psi)$ and
\[
 D(\Lambda)=\|2\mathbin{\times}\Lambda-\sqrt2\Lambda\|_2,
 \qquad H^2=D(\Lambda)^2/2.
\]
The identity follows by integration of the type-I step density. Pure local ancillas and the common amplification $h\mapsto I_9/9\otimes h$ preserve $H$.

Let $Q_j$ be Bob's label projections and $F_j$ Alice's copy flags. Set $d_j=\|(F_j-Q_j)\psi\|$, $N_L=\sum_jd_j^2$ and $N_F=\sum_jN_j$. The first links test exactly these mismatches, so their disjoint loss budgets give
\begin{equation}\label{eq:budgets}
 N_F+N_L\le243\eps.
\end{equation}
No commutation between the $F_j$ is assumed. For any ordering define the sequential instrument
\[
 T_r=F_r(1-F_{r-1})\cdots(1-F_1),\qquad
 T_{\rm fail}=(1-F_9)\cdots(1-F_1).
\]
Telescoping gives $\sum_rT_r^*T_r+T_{\rm fail}^*T_{\rm fail}=1$. Replace the rightmost remaining Alice factor on $\psi$ by its Bob counterpart, and commute that factor to the left. Orthogonality of the $Q_j$ gives
\[
 \|(T_r-Q_r)\psi\|\le\sum_{i\le r}d_i,\qquad
 \|T_{\rm fail}\psi\|\le\sum_id_i.
\]
Attach the sequential register on Alice and the exact label register on Bob. For $B=\sum_{j\in I}Q_j$, the target is $\sum_{j\in I}Q_j\psi\otimes|j,j\rangle$, a two-sided orthogonal sum. Squaring over register entries gives error squared at most
\[
 \sum_{r=1}^9\Bigl(\sum_{i\le r}d_i\Bigr)^2+
 \Bigl(\sum_id_i\Bigr)^2\le54N_L.
\]
Local isometries preserve Schmidt amplitudes, and singular-value contraction gives
\[
 \left\|\operatorname{Sch}(B\psi)
       -\bigcup_{j\in I}\operatorname{Sch}(Q_j\psi)\right\|_2
 \le\sqrt{54N_L}.
\]
Write $S_j=\operatorname{Sch}(Q_j\psi)$ and $U=\bigcup_jS_j$. Replacing $F_j$ by $Q_j$ in \eqref{eq:sector}, and splitting its complement, gives
\[
 \left\|\bigcup_{k\ne j}S_k-4\mathbin{\times}\sqrt2S_j\right\|_2
 \le C_F\sqrt{N_j}+(1+2\sqrt2)d_j+\sqrt{54N_L}.
\]
Concatenate and rearrange the nine comparisons. Each $S_k$ occurs eight times on the left and four times on the right. Hence
\[
 2D(U)\le C_F\sqrt{N_F}+(1+2\sqrt2+3\sqrt{54})\sqrt{N_L}.
\]
Also $\|\Lambda-U\|_2\le\sqrt{54N_L}$. Since $D$ is $2\sqrt2$-Lipschitz,
\[
 D(\Lambda)\le\tfrac{C_F}{2}\sqrt{N_F}+c_L\sqrt{N_L},\qquad
 c_L=\frac{1+2\sqrt2+3\sqrt{54}}2+2\sqrt{108},\quad c_L^2<1138.
\]
Cauchy--Schwarz with the disjoint budgets \eqref{eq:budgets} proves $H^2\le C_R\eps$ with the constant of Theorem~\ref{thm:price}. Only the first links entered; deleting the nine second links and renormalizing changes 243 to 234.

\paragraph{Semifinite implementation.}
After support reduction write the standard vector as $\zeta=h^{1/2}\in L^2(M,\Tr)$. Bob acts by right multiplication. Naimark is performed after compression: for a POVM column $W$ with $W^*W=e_{11}$, the block unitary
\[
 \begin{pmatrix}W&1-WW^*\\1-W^*W&W^*\end{pmatrix}
\]
implements its dilation after putting the unused second block into one outcome. Use unnormalized matrix traces, and embed $\zeta$ in one input corner. Thus its singular-value distribution is unchanged.

With $r$ left and $c$ right register entries, vectors are $r\times c$ matrices over $L^2(M)$, regarded as rectangular measurable operators in a linking algebra. Left column isometries and right row isometries preserve their nonzero singular distributions. Contractive multiplication decreases $\mu_t$; two-sided orthogonal sums add distribution functions. The generalized trace inequality \cite[Lemma 2.5 and Theorem 4.2]{FK} gives
\[
 \operatorname{Re}\Tr(y^*x)\le\int\mu_t(x)\mu_t(y)\,dt,
 \qquad \|\mu(x)-\mu(y)\|_2\le\|x-y\|_2.
\]
Truncation extends these statements to all $L^2$ elements. The extraction used in \eqref{eq:sector} is a finite column of bounded physical words and a finite group-register change of basis. Its common-invariant register projection is $|\Phi_4\rangle\langle\Phi_4|$; its range is $(I_4/2)\otimes j$, giving four copies of $\mu(j)/2$. The row/column rank-one product contributes $\mu(j)/(2\sqrt2)$. These are exactly the spectral operations in the square-root extraction and sector argument. Replacing Schmidt lists by rectangular singular distributions therefore transfers \eqref{eq:sector} with a universal finite constant, and every step of the sequential proof, without a matrix approximation assumption on $M$.

If Bob is the semifinite party, standardize his marginal, put the label on the left and the flags on the right, and take adjoints of rectangular coefficient operators. The identity $\mu(x)=\mu(x^*)$ preserves the argument. The game need not be symmetric under exchange of question sets. This proves the either-party form of the robust theorem.

\subsection{The exchange identity and the exact converse}
For a pure spatial vector, adjoining a Bell pair replaces its Schmidt list by $(2\mathbin{\times}\Lambda)/\sqrt2$. Spectral alignment attains the maximum overlap, so
\begin{equation}\label{eq:exchange}
 \inf_{U,V}\|(U\otimes V)\psi-\psi\otimes\Phi_2\|^2=H(h)^2.
\end{equation}
Kernels are padded by local pure ancillas. This is the usual Schmidt-overlap calculation; related pure-state conversion machinery appears in \cite{ZTG}. In a general semifinite representation singular-value contraction proves only the lower inequality for allowed maps. It supplies no operational upper bound for a smaller physical partner. In a standard full-commutant semifinite \emph{factor}, equality holds as an infimum: approximate both coefficient operators in $L^2$ by finite spectral step operators, split equal-trace spectral projections, and align their left and right supports by partial isometries. Equal-trace projections and their complements are equivalent in a finite factor; in an infinite factor first use finite-trace corners and then complete the infinite complements. The step approximants attain the rearranged overlap and $L^2$ convergence gives infimum $H^2$. Attainment for diffuse data is not asserted. This permits error at most $C_R\eps+\delta$ for every $\delta>0$ in that full-partner model. In a nonfactor, central labels can obstruct such alignment, so global rearrangement alone does not give this operational equality.

There is a profile-by-profile converse. On $\ell^2(\N)$ use $s_1|n\rangle=|2n-1\rangle$ and $s_2|n\rangle=|2n\rangle$. They are genuine Cuntz isometries, including the fixed point of $s_1$ at $n=1$. For a sorted probability density with amplitudes $\Lambda_n$, put it in the amplified marginal $I_9/9\otimes h$. The five-block construction has diagonal coefficients
\[
 (\Lambda_m,\Lambda_{\lceil m/2\rceil},\ldots,
                 \Lambda_{\lceil m/2\rceil})/3.
\]
On a block with coefficients $(a,b,b,b,b)$, rows, columns and calibrations lose zero and the total forbidden sign mass is $(\sqrt2a-b)^2$. Summing blocks and copies yields exactly
\begin{equation}\label{eq:converse}
 \eps=\frac1{243}\sum_m(\sqrt2\Lambda_m-\Lambda_{\lceil m/2\rceil})^2
      =\frac2{243}H(h)^2.
\end{equation}
Both links are exact. For unsorted lists the first equality remains true with their ordered defect. Finite support of size $n$ compresses the effects to genuine POVMs in dimension $9n$, preserving probabilities.

For $\eps_*(\sigma)$ optimized over type-I joint states and measurements with marginal $\sigma$, allowing Bob an unrestricted purification, we have
\[
 \frac{H(\sigma)^2}{C_R}\le\eps_*(\sigma),\qquad
 \eps_*(I_9/9\otimes h)\le\frac{H(h)^2}{121.5}.
\]
This is a two-sided comparison on specified amplified marginals. Equation~\eqref{eq:exchange} concerns pure states or mathematical purifications. It is false for the actual entanglement of an arbitrary mixed state: $\sum_ip_i|ii\rangle\langle ii|$ has the same marginal as its pure Schmidt counterpart, but branchwise product structure makes its squared Bell-exchange vector error at least $2-\sqrt2$ when its environment is untouched. Marginal entropy is likewise not mixed-state distillable entanglement.

\subsection{Finite templates, with their endpoints}
Let $c_0,\ldots,c_L\ge0$, $\sum c_l^2=1$, on the truncated binary-word construction, and write $w_l=c_l^2$. There are $9\cdot2^l$ equal eigenvalues $w_l/(9\cdot2^l)$ at level $l$. A parent block contributes the sign form just calculated; its summed interior mass is $2\sum_{l<L}(c_{l+1}-c_l)^2$. The empty-word flag and top-level magic deterministic completions contribute $c_0^2$ and $8c_L^2$. Therefore
\begin{equation}\label{eq:template}
 \eps(c)=\frac{2\sum_{l=0}^{L-1}(c_{l+1}-c_l)^2+c_0^2+8c_L^2}{243}.
\end{equation}
This includes both boundary penalties. Its resource identities are
\[
 S_I=\ln9+\mathcal H(w)+(\ln2)\mathbb E[l],
\]
\[
 D_{\mathrm{II}_1}=(\ln2)\mathbb E[L+1-l]-\mathcal H(w)
                   +\ln(1-2^{-(L+1)}),
\]
where $\mathcal H(w)=-\sum_lw_l\ln w_l$ and the finite-factor construction is a unital $M_{9(2^{L+1}-1)}$ embedding. For a bilateral realization, partition a $\mathrm{II}_\infty$ identity into levels with traces $9\cdot2^l$, and each level into nine equal atoms. Trace comparison identifies a source atom at level $l+1$ with each pair of target atoms at level $l$. Four such maps give the five-block system for each copy. The density is scalar $w_l/(9\cdot2^l)$ on each level, so the maps have exactly the coefficients used in the block loss calculation. Countable orthogonal sums give normal POVMs. Bilateral ladders in $\mathrm{II}_\infty$ thus have energy $E=\sum_{l\in\Z}(c_{l+1}-c_l)^2$ and loss $2E/243$, including every jump to a zero coefficient.

Let $a_1=2.338107410\ldots$ be the first positive Airy zero parameter and put
\[
 c_*=\frac{2a_1^{3/2}}{3\sqrt3},\qquad
 \kappa=(\ln2)c_*\sqrt{2/243}=0.0865330722761432\ldots.
\]
The normalized Dirichlet ground state of $-d^2/dx^2+x$ has kinetic energy $a_1/3$ and first moment $2a_1/3$, by scaling variation. Sampling its dilates on a half-line gives energy $a_1/(3r^2)+O(r^{-3})$, mean distance $2a_1r/3+O(1)$, and $\mathcal H(w)=\ln r+O(1)$. A shifted or vanishing boundary sample and a cutoff at $r$ times a power of $\ln r$ control the endpoint and Airy-tail errors. Inverting \eqref{eq:template} gives sufficient costs
\begin{align*}
 S_I(\eps)&\le\kappa\eps^{-1/2}+\tfrac12\ln(1/\eps)+O(1),\\
 D_{\mathrm{II}_1}(\eps)&\le\kappa\eps^{-1/2}-\tfrac12\ln(1/\eps)+O(1).
\end{align*}
These are abstract strategies, without a Hamiltonian interpretation.

For bilateral Gaussian $w_l\propto e^{-l^2/(2\sigma^2)}$, the shift energy is $2[1-e^{-1/(8\sigma^2)}]$ up to exponentially small lattice errors. Thus $\eps\sim1/(486\sigma^2)$. Symmetry makes $\operatorname{Cov}(l,\ln w_l)=0$ and $\Var(\ln w_l)\to1/2$, giving
\[
 \Var_\phi(\ln h)=(\ln2)^2\sigma^2+\tfrac12+o(1),\qquad
 \eps\Var_\phi(\ln h)\longrightarrow\frac{(\ln2)^2}{486}.
\]
For arbitrary profiles the covariance term must be retained; the difference between log-density and level variances need not be uniformly bounded.

\subsection{Measure bounds and the sharper wall estimates}
The density of $X$ is $f(s)=(\ln2)2^sp(2^s)$; with $g=\sqrt f$,
\[
 H^2=\int_\R(g(s+1)-g(s))^2\,ds.
\]
Fubini with $s=u+k$, $0\le u<1$, fibres the mass, energy and first moment into discrete ladders; the offset $u$ stays in the moment. A nonzero probability cannot have $H=0$.

For any unit-interval masses $m_j$, Cauchy--Schwarz gives $\sum_j|m_j-m_{j-1}|\le2H$. A summable nonnegative sequence has variation at least $2\max m_j$, so every unit interval has mass at most $H$. In type I, $p=\lambda_{\max}$ on $(0,1)$, whence $\lambda_{\max}\le2H$. Covering $[0,x]$ by at most $x+1$ unit intervals and integrating the surviving tail gives the valid all-$H$ bound
\begin{equation}\label{eq:elementary}
 S(h)\ge\frac{(\ln2)(1-3H)_+^2}{2H}-1.
\end{equation}
Here $\mathbb E[X_-]=\lambda_{\max}/\ln2$ and $tp(t)\le1$ were used. The positive part is essential.

For normalized finite trace, decreasing $p$ on $(0,1)$ dominates the uniform cumulative distribution, so $m=\mathbb E[-\ln T]\ge1$. Relative entropy against $t^{-\theta}$ gives $D\ge\theta m+\ln(1-\theta)$, $\theta<1$, and optimization yields
\begin{equation}\label{eq:deficit}
 D\ge m-1-\ln m.
\end{equation}
A lower estimate inserted for $m$ must be replaced by its maximum with 1. Infinite moments follow by the variational form before optimization.

To get sharp leading wall constants, let $b_n$, $n\ge0$, have mass $M$, energy $D=\sum(b_{n+1}-b_n)^2$ including $b_{-1}=0$, and moment $A=\sum nb_n^2$. Piecewise-linear interpolation has squared norm $M-D/6$, derivative energy $D$, and shifted first moment at most $A+3M/2$. The continuum Airy variational inequality therefore gives, for $\theta>0$,
\[
 \theta D+A\ge
 M\frac{a_1\theta^{1/3}-3/2}{1+a_1\theta^{-2/3}/6}.
\]
Truncation includes infinite moment cases. This is a variational proof of the discrete estimate, not a fit to numerical eigenvalues.

In type I integration over $(0,1/2)$ in the definition of $H$ gives $\lambda_{\max}\le aH^2$, $a=2/(\sqrt2-1)^2$. Cut $g$ below $-K$, $K=\lceil2\log_2(1/H)\rceil$. The removed mass is $\delta\le aH^4$, remaining mass $M=1-\delta$, and new energy $E\le H^2+\sqrt2\delta$. The removed moment is exactly $-\delta(K+1/\ln2)$. Fibre the Airy inequality and optimize $\theta=(a_1M/(3E))^{3/2}$ to obtain
\[
 \mathbb E[X]\ge\frac{c_*M^{3/2}}{\sqrt E}-\frac{3M}{2}
 -\frac{3c_*}{4}\sqrt{ME}-K-\frac\delta{\ln2}.
\]
For $0<H\le0.05$, elementary bounds on these explicit remainders yield
\begin{equation}\label{eq:wall}
 \mathbb E[X]\ge\frac{c_*}H-2\log_2(1/H)-3.5,\qquad
 S(h)\ge(\ln2)\mathbb E[X].
\end{equation}
Indeed the additional remainder is at most
$c_*(\sqrt2a/2+3/4)H+c_*(3a/2+3\sqrt2a/8)H^3+aH^4/\ln2<1$.

For rank $q$, each cutoff fibre has at most $\log_2q+K+1$ sites. The Dirichlet box eigenvalue $2[1-\cos(\pi/(n+1))]$ yields
\[
 \ln q\ge\frac{(\ln2)\pi}H-O(\ln(1/H)).
\]
For normalized $\mathrm{II}_1$ trace the upper wall at $T=1$ is exact. Reversing ladders gives
\[
 \mathbb E[-X]\ge\frac{c_*}H-\frac32-\frac{3c_*}{4}H.
\]
Combine it with \eqref{eq:deficit}, using $m\ge1$, to get $D\ge(\ln2)c_*/H-\ln(1/H)-O(1)$.

For variance, the translated laws $f(s)$ and $f(s+1)$ have means separated by one. Centering at their midpoint gives
\[
 1\le H\sqrt{4\Var(X)+1},\qquad
 \Var(X)\ge\frac{1-H^2}{4H^2}.
\]
Infinite variance satisfies the extended-real statement.

Finally put $K_{\rm mod}=-\ln p(T)$ and $W=K_{\rm mod}-\ln T\ge0$. With $V=\Var(X)<\infty$, a length-$2L$ interval and Chebyshev give
\[
 \Pr(W>w)\le V/L^2+2L(\ln2)e^{-w}
 \le3(\ln2)^{2/3}V^{1/3}e^{-2w/3}.
\]
For $w_0=\tfrac12\ln V+\tfrac32\ln3+\ln(\ln2)$ in the small-$H$ range, tail integration gives $\mathbb E[W^2]\le w_0^2+3w_0+9/2$. The triangle inequality for standard deviations gives the safe bound
\[
 \Var(K_{\rm mod})\ge
 \left[(\ln2)\sqrt V-\sqrt{w_0^2+3w_0+9/2}\right]_+^2.
\]
For $H\le0.05$ the bracket is positive and increasing for $V\ge(1-H^2)/(4H^2)$, so
\[
 \Var(K_{\rm mod})\ge\frac{(\ln2)^2}{4H^2}
                    [1-O(H\ln(1/H))].
\]
To cover $V=\infty$, let $Y=\ln T$. Monotonicity gives $Y\le K_{\rm mod}$ and, for any real $b$ and $t>0$,
\[
 \Pr(Y<b-t)\le\Pr(K_{\rm mod}<b-t/2)+e^{-t/2}.
\]
Integrating the two tails yields $\mathbb E[(Y-b)^2]\le4\mathbb E[(K_{\rm mod}-b)^2]+8$. Finite modular variance would imply finite $V$, proving the remaining case.

\subsection{Sorted trace position and matching exponents}
An unsorted bilateral ladder density is $r(t)=w_l/(9\cdot2^l)$ on $[9\cdot2^l,9\cdot2^{l+1})$. Its ordered defect is exactly $E$. The actual singular function is $r^*$. Rearrangement increases affinity, since intersections of superlevel sets are maximized by initial intervals, and $(2r(2t))^*=2r^*(2t)$. Thus
\[
 H_{\rm sorted}^2\le E=121.5\eps,
\]
with equality only under an appropriate monotone level ordering. Applying the variance inequality gives the sharp template result
\[
 \lim_{\eps\downarrow0}\eps\inf_{\substack{\text{bilateral templates}\\\ell\le\eps}}\Var(\ln T)
       =\frac{(\ln2)^2}{486}.
\]
For the upper direction use a broad Gaussian and, below $l=-\sigma^2\ln2/2$, continue its tail geometrically with adjacent ratio $e^{(\ln2)/2}$. Then $w_l/2^l$ is decreasing; normalization, energy and moments change only exponentially little. Conditional on a level, $T=9\cdot2^lU$ with $U$ uniform on $[1,2]$, giving $\Var(\log_2T)=\sigma^2+O(1)$. This proof concerns sorted trace position. The Gaussian log-density upper bound is a separate statement.

Combining the robust law and measure bounds with the Airy, sine and Gaussian constructions gives the following leading terms. Negative logarithmic remainders accompany the first three lower bounds; the last lower term has relative $O(\sqrt\eps\ln(1/\eps))$ remainder, with fixed-$C_R$ constants.
\[
\begin{array}{l|cc}
\text{resource}&\text{lower leading term}&\text{attainable upper leading term}\\\hline
 S_I&(\ln2)c_*/\sqrt{C_R\eps}&(\ln2)c_*/\sqrt{121.5\eps}\\
 \ln d&(\ln2)\pi/\sqrt{C_R\eps}&(\ln2)\pi/\sqrt{121.5\eps}\\
 D_{\mathrm{II}_1}&(\ln2)c_*/\sqrt{C_R\eps}&(\ln2)c_*/\sqrt{121.5\eps}\\
 \Var(\log_2T)&1/(4C_R\eps)&1/(486\eps)\\
 \Var_\phi(-\ln h)&(\ln2)^2/(4C_R\eps)&(\ln2)^2/(486\eps)
\end{array}
\]
The trace-position lower inequality is exactly $1/(4C_R\eps)-1/4$. A sine profile with zero endpoints on a length-$L$ ladder has energy $2[1-\cos(\pi/L)]$ and dimension logarithm $L\ln2+O(1)$. Together with the rank lower bound this gives
\[
 \frac{((\ln2)\pi)^2}{C_R(\ln d)^2}(1-o(1))
 \le1-\omega_d(\GIII)
 \le\frac{((\ln2)\pi)^2}{121.5(\ln d)^2}(1+o(1)).
\]
This proves Theorem~\ref{thm:price}, for POVMs and mixed strategies as well. The sine mechanism is prior methodology \cite{LTW,Precision}, and exponential game witnesses precede this result \cite{Coladangelo,Explicit}. The authors' explicit-game companion already asserts a different game's law $\exp(\eps^{-1+o(1)})$; its $\Theta(\eps^{-1/2})$ law for another game concerns \emph{dimension}, not its logarithm. The present result determines this game's $\exp(\Theta(\eps^{-1/2}))$ law. It does not improve a universal previous-best exponent of $1/8$.

These are asymptotic bounds with a universal finite $C_R$. The sufficient small-$H$ range is $\eps\le0.05^2/C_R$. No validated numerical constant is given; the sharp constant is conjectural. The proposed sharp robustness coefficient $C_R=121.5$ remains conjectural; the exact converse proves its necessity as a lower candidate, not its universal validity.

\subsection{Deficit divergence without a quantitative estimate}
There is also a direct compactness proof for normalized finite algebras. Suppose losses tend to zero with $D_{\tau_n}(h_n)\le C$. Support reduction and trace normalization change the deficit to $D+\ln\tau_n(p)\le D$. Standardize both sides. Since $x\ln x\ge-1/e$,
\[
 \tau_n(h_n1_{(K,\infty)}(h_n))\le\frac{C+1/e}{\ln K}=:t_K\longrightarrow0.
\]
The truncated vectors $\zeta_{n,K}=\sqrt{\min(h_n,K)}$ have bounded operator norm and squared $L^2$ error at most $t_K$. Their classes in the tracial ultraproduct converge in $L^2$ to a unit vector. All coordinate effects give bounded POVMs there. A correlation changes by at most $2\sqrt{t_K}$ upon truncation, so the finite list of limiting correlations is perfect. The ultraproduct is finite and its vector state normal, contradicting semifinite exclusion. Consequently bounded deficit forces a positive loss gap, uniformly over these normalized finite-algebra strategies. It does not bound the deficit of every state in a fixed finite factor.

\section{Type versus price}\label{sec:phase}
For finite rational games, commuting value one is attained; bounded dimension cost is equivalent to a finite-dimensional perfect strategy, and a normal type-I perfect host precludes a quantum value gap. A matrix rank/dilation-defect inequality with sharp leading constant $\pi$ converts a proved power bound on the defect into an exponential dimension lower bound. It does not characterize all type-III-only games. Published quantitative games and $\GIII$ give polynomial, exponential and value-gap examples of minimum type II, and exponential and value-gap examples of minimum type III. Every $\Delta^0_2$ exponent in $[1/2,1]$ is realized by a type-III-only game; every positive finite existing exponent is $\Delta^0_2$. The game $\GIII$ attains the exponent floor within the first-order certification class. The perfect-commuting language is $\Pi^0_1$-complete; type-III-only hosting is d.r.e.-hard, and its intersection with quantum value one is $\Pi^0_2$-hard as an unpromised language. Exact arithmetical classification, the promised-input strengthening and the universal type-to-price law remain open here. The precise decoupling is that minimum type II does not select a unique price class, and exponential price does not determine whether the minimum type is II or III.

Write $\omega_q=\sup_d\omega_d$ and $\omega_{qc}$ for the commuting value. Among hosted games let the \emph{minimum host class} be the least of
\[
 \text{finite matrices}<\mathrm I_\infty<\mathrm{II}<\mathrm{III}.
\]
Here II includes $\mathrm{II}_1$ and $\mathrm{II}_\infty$; ``finite matrices'' does not include finite type-II algebras. The last class means that every perfect normal host is nonsemifinite. Set $D_G(\eps)=\infty$ if the specified loss is infeasible.

\subsection{Standard attainment and truncation facts}
\begin{proposition}[Standard compactness and truncation]\label{prop:attainment}
Every finite game of commuting value one has a perfect normal strategy on a factor with separable predual. Its dimension cost is bounded as $\eps\downarrow0$ exactly when a matrix factor hosts it. If it has a normal type-I perfect host with pure Schmidt weights $(\lambda_k)$, then
\[
 D_G(\eps)\le\min\{d: \textstyle\sum_{k>d}\lambda_k\le\eps/2\}<\infty.
\]
In particular a type-I perfect host implies $\omega_q=1$.
\end{proposition}
\begin{proof}
These are standard recoveries, not new compactness results \cite{NPA,MNY}. The compact state space of the universal measurement algebras' maximal tensor product maps continuously onto $C_{qc}$, so the payoff attains its maximum. Separable GNS and central decomposition of Alice's generated algebra give factor fibers. Bob commutes with its center and is decomposable; finitely many nonnegative losing terms vanish on almost every fiber. The vector states are normal.

Fixed-dimensional states and POVMs form a compact set, proving the bounded-cost equivalence. A normal type-I representation is an amplification, hence spatial. From a perfect mixed state select a pure zero-loss component and write $\psi$ in Schmidt form. If $t_d=\sum_{k>d}\lambda_k$, its normalized truncation obeys, for the positive loss contraction $L$ with $L\psi=0$,
\[
 \langle\psi_d,L\psi_d\rangle\le\|\psi_d-\psi\|^2
       =2-2\sqrt{1-t_d}\le2t_d.
\]
Compress both POVMs to the corresponding Schmidt subspaces.
\end{proof}

\subsection{A quantitative implication from dilation defect}
For $r>1$ define
\[
 H_r(h)^2=\int_0^\infty\left(\sqrt{p(t)}-\sqrt{rp(rt)}\right)^2dt,
 \qquad p=\mu(h),\quad\int p=1.
\]
At each \emph{feasible} loss $\eps$, let $\delta_{G,r}(\eps)$ be the supremum of $H_r(\rho_A)$ over finite-dimensional strategies of loss at most $\eps$, using the canonical counting trace on matrices. Let $\delta^{\rm sf}_{G,r}$ use all normal semifinite strategies where this set is nonempty. No positive modulus is assigned to an empty strategy set. A power certificate means $\delta_{G,r}(\eps)\le C\eps^\gamma$; first-order certification means $\gamma=1/2$.

\begin{theorem}[Matrix rank and dilation defect]\label{thm:rank-defect}
For each $r>1$ there is $c_r>0$ such that a matrix density of rank $R$ and defect $0<H=H_r\le c_r$ satisfies
\begin{equation}\label{eq:rank-defect}
 \log_r R\ge\frac\pi H-2\log_r(1/H)-3.
\end{equation}
One may take $c_2=.05$. The leading coefficient $\pi$ is sharp for matrix densities. Consequently, at feasible losses with $0<\delta_{G,r}(\eps)\le c_r$, the same bound holds for $\log_r D_G(\eps)$ with $H$ replaced by $\delta_{G,r}(\eps)$. A proved power bound on the defect gives an exponential dimension lower bound. Vanishing $\delta^{\rm sf}_{G,r}(\eps)$ as $\eps\downarrow0$ implies semifinite exclusion.
\end{theorem}
\begin{proof}
Put $f(s)=(\ln r)r^sp(r^s)$ and $g=\sqrt f$. Then $\|g\|_2=1$ and $H^2=\int|g(s+1)-g(s)|^2ds$. The matrix bottom plateau $p=\lambda_1$ on $(0,1]$ gives
\[
 \lambda_1\le a_rH^2,\qquad a_r=\frac r{(\sqrt r-1)^2}.
\]
Cut $g$ below the nonintegral point $s_0=2\log_rH$. For $H<r^{-1/2}$ the crossing edges lie in the exponential tail. The removed mass and the cut function's energy satisfy
\[
 d=\lambda_1r^{s_0}\le a_rH^4,\quad M=1-d,\qquad
 E\le H^2+2(\sqrt r-1)d.
\]
Indeed deleting edges wholly below $s_0$ saves $(\sqrt r-1)^2d$, while replacing crossing edges by zero-boundary edges costs $(r-1)d$; their difference is $2(\sqrt r-1)d$. The retained interval has length $B=\log_rR+2\log_r(1/H)$. Fiber modulo one. Each fiber has at most $B+1$ sites, so the exact Dirichlet ground energy gives $E\ge4\sin^2(\pi/(2(B+2)))M$. Therefore
\[
 \log_rR\ge\frac\pi{2\arcsin z_r(H)}-2\log_r(1/H)-2,
 \quad z_r(H)=\frac H2\sqrt{\frac{1+b_rH^2}{1-a_rH^4}},
 \quad b_r=2(\sqrt r-1)a_r.
\]
For $r=2$ and $H\le.05$, the square-root factor is at most $1+4.9H^2$. Using
\[
 \frac{\arcsin z}{z}\le1+\frac{z^2}{6(1-z^2)^{3/2}}
\]
gives $2\arcsin z_2(H)\le H(1+5H^2)$, hence $\pi/(2\arcsin z_2(H))\ge\pi/H-5\pi H>\pi/H-1$. This proves the stated remainder. For fixed $r$, the difference from $\pi/H$ is $O_r(H)$; decrease $c_r$ to bound it by one and make the right side of \eqref{eq:rank-defect} decreasing in $H$. Substitution of $H\le\delta$ proves the game inequality.

For sharpness fix $a>2/\ln r$. On $[a,L]$ use normalized sine amplitude $g(s)\propto\sin(\pi s/L)$, continue below $a$ by $g(a)r^{(s-a)/2}$ and set it zero above $L$. The corresponding $p$ is constant at the bottom and decreasing above, since its logarithmic derivative is at most $2/a-\ln r<0$. The cap contributes $O_r(L^{-3})$ to mass and energy, while the bulk shift energy is $\pi^2/L^2+O_r(L^{-3})$. Average $p$ on unit trace intervals. The squared $L^2$ error of the square-root densities is $O_r(L^{-3})$, by summing cell oscillations and bounding $\int|(\sqrt p)'|^2$. The dilation difference operator has norm at most two. Thus $H=\pi/L+o(L^{-1})$, $R=\lceil r^L\rceil$, and $(\log_rR)H\to\pi$.

Finally $H_r(h)>0$ for every normal semifinite density: zero would give a nonzero finite translation-invariant probability law $f(s)ds$. A perfect strategy would be counted at every positive loss and contradict a vanishing semifinite modulus.
\end{proof}
The matrix plateau is essential; diffuse densities need not obey this rank argument. In the displayed $\GIII$ proof, Theorem~\ref{thm:price} supplies $H_2\le\sqrt{C_R\eps}$, so
\[
 \log_2D_{\GIII}(\eps)\ge\frac\pi{\sqrt{C_R\eps}}
             -O(\log(1/\eps)).
\]
The sine construction gives the matching exponent and upper coefficient $\pi\sqrt{2/243}$. Robustness supplies the rate in this proof; exclusion and fixed-dimensional compactness alone give divergence. No impossibility theorem about other proofs is asserted. Nor is there a converse to the defect implication: flat marginals satisfy $H_r(I_d/d)^2=2(1-r^{-1/2})$, independent of $d$, even in the exponentially expensive examples below.

\subsection{Composition, type-II examples and programmable exponents}
Define $G_1\wedge G_2$ by product questions and answers, winning exactly when both components win. For $G_1\vee G_2$, every product-question answer carries a mode $b\in\{1,2\}$ and an answer for that component. With probability $1/2$ play product questions and test agreement of modes and the selected game; with probability $1/4$ each, send a mode-only question $M$ to one party and a uniform product question to the other, testing just mode agreement. Remove questions of zero marginal weight first. The two mode stars and a supported product edge make the mode graph connected.

\begin{proposition}[AND/OR hosting and prices]\label{prop:boolean}
If $M$ hosts $G_1\wedge G_2$, it hosts both components. Hosts $M_1,M_2$ give a host $M_1\bar\otimes M_2$ for the product. For every factor $M$,
\[
 M\text{ hosts }G_1\vee G_2\iff
 M\text{ hosts }G_1\text{ or }G_2.
\]
When the relevant hosts exist, minimum host classes take maximum for AND and minimum for OR. Dimension costs satisfy
\begin{align*}
 \max_iD_i(\eps)&\le D_{G_1\wedge G_2}(\eps)
                  \le D_1(\eps/2)D_2(\eps/2),\\
 \min_iD_i(C\eps)&\le cD_{G_1\vee G_2}(\eps),\qquad
 D_{G_1\vee G_2}(\eps)\le\min_iD_i(2\eps)
\end{align*}
for fixed positive $c,C$ and sufficiently small feasible $\eps$.
\end{proposition}
\begin{proof}
For AND, fix unused questions and coarse-grain their answers. At positive loss select a pair by averaging. Tensor strategies and the union bound give the upper bound; tensor products give the stated maximum of coarse host classes.

For perfect OR, support reduction and Lemma~\ref{lem:classes}, propagated through the connected mode graph, give one centralizer projection $z$ and its commutant mirror, not a central projection. If $\phi(z)>0$, condition on $z\xi$ and restrict to $zMz$; coarse-graining the unused component gives a perfect $G_1$ strategy. Otherwise use $1-z$ and $G_2$. A corner host extends to $M$: its standard vector lies in $zJzJL^2(M)$, left and right effects sum to $z$ and $JzJ$, and completing on their complements preserves commutation and correlations. Constant mode gives the converse.

For approximate OR select a pure component of no greater loss, then use common local Naimark dilations of constant dimension cost. All mode and fine effects are projections. The smallest supported question weight $\mu_0>0$ bounds each mode-disagreement norm by $\sqrt{\eps/\mu_0}$. Paths in the mode graph give vector distances $\eta=K\sqrt\eps$. Choose a mode at $M$ with Alice weight at least $1/2$, with mode projections $P,Q$ on the two parties, and put $\theta=PQ\psi$. Then
\[
 \|\theta\|^2\ge\tfrac12-\eta^2,\qquad
 \|\theta-A_x^bB_y^b\psi\|\le2\eta.
\]
The losing fine projections in that mode sum to a projection $L_{xy}^b$ supported on $A_x^bB_y^b$. Hence
\[
 \|L_{xy}^b\theta\|\le\sqrt{\eps/\mu_0}+2\eta.
\]
On normalized $\theta$ this is loss $O(\eps)$. Compress fine measurements to $P,Q$. Their missing effects $P(1-A_x^b)P$ and $Q(1-B_y^b)Q$ have expectation $O(\eta^2)$ on $\theta$; assign each to one allowed answer. This completes the POVMs at extra loss $O(\eps)$. Fix and coarse-grain unused questions by averaging to extract $G_b$ in dimension at most $c$ times the original. Constant-mode play gives the OR upper bound, since only its product-question half can lose. These are inequalities, not literal maximum/minimum identities of price functions.
\end{proof}

\begin{proposition}[An exponential game of minimum type II]\label{prop:ii-exp}
The authors' explicit linear-system game $G_{\rm ex}$ has a $\mathrm{II}_1$ factor host and no normal type-I host, while $\omega_q(G_{\rm ex})=1$ and $\log D_{G_{\rm ex}}(\eps)=\eps^{-1+o(1)}$ \cite{Explicit}.
\end{proposition}
\begin{proof}
Commuting perfection follows from Proposition~\ref{prop:attainment}. The solution-group criterion \cite[Theorems 4--5]{CLS} makes its central involution $J$ nontrivial. In $L(\Gamma)$, $p=(1-J)/2$ has trace $1/2$. Left/right regular play on $L(\Gamma)p$, with vector $\sqrt2p\delta_e$, is perfect and tracial. A nonzero type-I part of this finite algebra would disintegrate into matrices; one fiber would detect $J=-1$ and yield finite perfect play, which this game excludes. Thus its factor decomposition supplies a $\mathrm{II}_1$ host. Perfect spatial linear-system play in any dimension implies finite perfect play \cite[Theorem 3.2]{SloNonclosure}; normal type-I hosting is spatial, so no such host exists.
\end{proof}
The regular algebra need not be a factor. This argument identifies neither arbitrary ambient type-III hosts nor approximation properties of its canonical regular correlation. The polynomial type-II example of Slofstra--Vidick \cite{SV} has its matching $D=\Theta(\eps^{-1/2})$ sharpening in the authors' explicit-game paper \cite{Explicit}; the original bounds in \cite{SV} have unequal exponents. The general constraint-system examples in \cite{Uncomputable} also have minimum type II: shared variable mirrors make the generated state tracial; a matrix summand gives finite perfect play. In a spatial perfect strategy the trace-class marginal commutes with these generators, and any positive eigenvalue gives a finite invariant eigenspace and finite perfect play. This rules out type I when finite perfection is absent. The earlier type-II tracial-limit mechanism and dimension divergence are discussed in \cite{Precision}.

\begin{theorem}[Positive finite exponent spectrum]\label{thm:exponents}
Let $\mathcal E_{\rm III}^+$ comprise the positive finite existing limits $\beta(G)=\lim_{\eps\downarrow0}\log\log D_G(\eps)/\log(1/\eps)$ for rational type-III-only games with $\omega_q=1$. Then
\[
 \Delta^0_2\cap[1/2,1]\ \subseteq\ \mathcal E_{\rm III}^+
       \ \subseteq\ \Delta^0_2\cap(0,\infty).
\]
The upper restriction holds for every rational game with quantum value one and a positive finite existing exponent. Within the first-order certification class the lower exponent floor is $1/2$, attained by $\GIII$.
\end{theorem}
\begin{proof}
The authors' realization theorem \cite[Theorems 1.1 and 1.3]{Uncomputable} supplies fixed rational constraint games $G_\beta$ with $\log D_{G_\beta}=\eps^{-\beta+o(1)}$ for every $\Delta^0_2$ real $0<\beta\le1$. Put $G=\GIII\wedge G_\beta$. Tensor hosts and approximants give type-III-only hosting and quantum value one. The correct sandwich is
\[
 \max\{\log D_{\GIII}(\eps),\log D_{G_\beta}(\eps)\}
 \le\log D_G(\eps)
 \le\log D_{\GIII}(\eps/2)+\log D_{G_\beta}(\eps/2).
\]
Thus $\beta(G)=\max\{1/2,\beta\}$, proving realization, including noncomputable exponents. For the converse, fixed-dimensional feasibility at rational loss $\le\eps$ is a compact rational semialgebraic decision problem. Quantifier elimination and increasing dimension search compute the actual $D_G(\eps)$ under $\omega_q=1$, including equality at the threshold. Computable rational approximations to $\log(1+\log D_G(2^{-n}))/(n\ln2)$ have the same finite limit, which is $\Delta^0_2$ by the limit lemma; compare \cite[Proposition P6.1]{Uncomputable}. The first-order floor follows from Theorem~\ref{thm:rank-defect}. This proves no universal positive floor, ceiling at one, or existence of a limit for arbitrary games.
\end{proof}

\subsection{Complexity and the settled cells}
For rational encoded games put $\mathrm{PERF}=\{G: \omega_{qc}(G)=1\}$, let $\mathrm{SEMI}$ mean existence of a normal semifinite factor host, and put $\mathrm{III\!\text{-}\!ONLY}=\mathrm{PERF}\setminus\mathrm{SEMI}$ and $Z=\{G: \omega_q(G)=1\}$.
\begin{theorem}[Unpromised complexity lower bounds]\label{thm:type-complexity}
$\mathrm{PERF}$ is $\Pi^0_1$-complete (a prior result). The language $\mathrm{III\!\text{-}\!ONLY}$ is d.r.e.-hard under computable many-one reductions, and $\mathrm{III\!\text{-}\!ONLY}\cap Z$ is $\Pi^0_2$-hard as an unpromised language. The exact arithmetical placement of type-III-only hosting remains open here.
\end{theorem}
\begin{proof}
Mousavi--Nezhadi--Yuen \cite[Theorems 6.4 and 6.9]{MNY} give $\Pi^0_1$-completeness of exact commuting value one and $\Pi^0_2$-completeness of exact quantum value one. The first also follows from the bounded NPA moment hierarchy \cite{NPA} and Slofstra's effective embedding \cite[Corollary 3.3]{SloEmbedding}. Finite-level non-strict perfect moment feasibility is decidable and holds at every level exactly for commuting perfection.

Uniformly choose linear-system games $L_X$ perfect exactly when machine $X$ never halts. Their perfect strategies have regular tracial hosts, so $\mathrm{SEMI}=\mathrm{PERF}$ on these games. With $H=\GIII$, set $F(P,Q)=(H\wedge L_P)\vee L_Q$. The AND branch always excludes semifinite hosts and is hosted exactly when $P$ never halts; the OR alternative has a semifinite host exactly when $Q$ never halts. Hence
\[
 F(P,Q)\in\mathrm{III\!\text{-}\!ONLY}
 \iff Q\text{ halts and }P\text{ never halts}.
\]
Choosing the two halting sets independently reduces every difference of r.e. sets. In particular this language is neither r.e. nor co-r.e. Finally $L\mapsto H\wedge L$ reduces $Z$ to $\mathrm{III\!\text{-}\!ONLY}\cap Z$: product value one is equivalent to $\omega_q(L)=1$, and on those instances a tensor host exists. No-instances need not map into $Z$, so this is not hardness under a value-one promise.
\end{proof}
Lin's $\mathsf{MIP}^{co}=\mathsf{coRE}$ \cite[Theorem 1.1]{Lin} is the prior gapped commuting counterpart; Ji--Natarajan--Vidick--Wright--Yuen \cite{JNVWY} establish the gapped quantum complexity and value-separation context. Neither identifies the arithmetical complexity of semifinite hosting.

For regular profiles write B for bounded cost; L for unbounded subpolynomial cost; P for polynomial cost; SPE for superpolynomial cost with subpower $\log D$; E for an exponential lower bound $\log D\ge c\eps^{-\gamma}$ for some $c,\gamma>0$; and gap for $\omega_q<1$. These labels are not asserted to partition all oscillatory profiles. A plus denotes a construction, a minus an impossibility, and a question mark an open cell.
\begin{center}
\begin{tabular}{lcccccc}
\toprule
Minimum host & B & L & P & SPE & E & gap\\
\midrule
Finite matrices & $+$ & $-$ & $-$ & $-$ & $-$ & $-$\\
$\mathrm I_\infty$ & $-$ & $+$ & $?$ & $?$ & $?$ & $-$\\
II & $-$ & $?$ & $+$ & $?$ & $+$ & $+$\\
III only & $-$ & $?$ & $?$ & $?$ & $+$ & $+$\\
\bottomrule
\end{tabular}
\end{center}
There are sixteen settled cells and eight open cells. Magic square fills finite/B. The authors' base game has spatial perfection, no finite perfection and $D=10\log_2(1/\eps)+O(\log\log(1/\eps))$ \cite{Base}, filling $\mathrm I_\infty$/L. Proposition~\ref{prop:attainment} excludes every other finite row cell, all nonfinite/B cells and type-I/gap. The II/P and II/E constructions are those above. Gao--Yu--Zhi's explicit binary linear-system game $G_{\rm gap}$ satisfies $\omega_q<1=\omega_{qc}$ \cite[Theorems 1.1--1.2]{GYZ}; its regular tracial sector and type-I exclusion fill II/gap. The games $\GIII$ and $\GIII\wedge G_{\rm gap}$ fill III/E and III/gap. Exponential dimension witnesses already occur in Coladangelo's two-player construction \cite{Coladangelo} and Ji--Leung--Vidick's three-player coherent-embezzlement construction \cite{JLV}; these are prior quantitative phenomena with different task hypotheses.

Only games with a proved defect power certificate inherit the corresponding spectral exclusions; a necessary ladder alone supplies no price theorem. The settled cells show exactly this decoupling: minimum type II does not select a unique price class, and exponential price does not determine whether the minimum type is II or III. They do not establish full independence or settle whether every type-III-only value-one game has exponential price.

\section{Conditional applications and remaining questions}\label{sec:scope}
\begin{corollary}[Specified semifinite gravitational models]\label{cor:gravity}
Suppose the actual measurement algebras obey $A\subseteq N\subseteq B'$, where $N$ is one of the specified semifinite perturbative algebras below, and the strategy's state is normal on $N$. They cannot attain perfect $\GIII$ play. A fixed-background quantum-field-theory algebra already known to be a type-$\mathrm{III}_1$ factor with separable predual has perfect prepared-state hosting and the two-game universal-state response property with its standard commutant.
\end{corollary}
\begin{proof}
Apply the separator exclusion of Section~\ref{sec:games} to $N$. Apply Theorems~\ref{thm:host} and \ref{thm:two} to the stipulated type-$\mathrm{III}_1$ algebra. A physical complementary algebra can replace the full commutant only when duality or another justified identification supplies it.
\end{proof}
The source constructions, with their hypotheses retained, are these:
\begin{itemize}
\item Chandrasekaran--Longo--Penington--Witten's de Sitter observer patch algebra is type $\mathrm{II}_1$ after the specified finite compression \cite[Sections 2.4 and 4]{CLPW}. This does not turn a compression map into a homomorphism on the original type-III algebra. A finite algebra cannot contain a type-III von Neumann subalgebra: its faithful finite trace restricts to one on any subalgebra.
\item Witten's perturbative black-hole crossed product is type $\mathrm{II}_\infty$ \cite[Section 3]{Witten}; Chandrasekaran--Penington--Witten give the relevant type-II microcanonical exterior construction \cite[Sections 2--4]{CPW}. Retain their ensemble, region and perturbative qualifications. The type-III algebra can remain a subalgebra of a $\mathrm{II}_\infty$ core, but Bob commuting with the whole core must also commute with its additional generator. Its original undressed commutant cannot simply be supplied as the partner.
\item Jensen--Sorce--Speranza's compact observer and boundary constructions yield, respectively, $\mathrm{II}_1$ and $\mathrm{II}_\infty$ under their stated assumptions \cite[Sections 2.1 and 5]{JSS}. The required kinematical type, constraint and modular/geometric-flow state remain hypotheses; the bounded-observer conclusion also retains their lower-energy-bound assumption; assumed state existence is not proved by the present game.
\item Kudler-Flam--Leutheusser--Satishchandran's explicitly typed full-algebra cases are semifinite, with their stationary-state and decay hypotheses \cite[Sections 6--8]{KFLS}. These include the Kerr $\mathrm{II}_\infty\bar\otimes\mathrm{I}_\infty$ description and the two $\mathrm{II}_\infty$ horizon factors in Schwarzschild--de Sitter. Mere containment of a type-II subfactor would not establish semifiniteness of a larger algebra.
\end{itemize}
The broader task-level observation that gravity's proposed change to type II removes normal embezzling resources was already made by van Luijk--Stottmeister--Werner--Wilming \cite[Corollary 54 and Section 2, version 4]{VLSWW}. Their fixed-resource universal-embezzlement task must be distinguished from finite-error protocols whose resource state changes with the error. Our corollary is its particular finite classical-game counterpart.

Each individual normal semifinite strategy has positive loss, while the unrestricted infimum over finite-dimensional strategies is zero. A positive resource-capped gap requires the cap itself as an additional premise. No numerical de Sitter floor, Hamiltonian or preparation cost follows here. Nor does excluding the listed models exclude every crossed product: a trivial-action product of a type-III host with an abelian group algebra still carries a perfect strategy. The conclusions concern the specified normal semifinite separators.

\subsection{Open problems}
Theorem~\ref{thm:rationals} supplies every positive rational, while Theorem~\ref{thm:orthogonal} determines the orthogonal-pair single-bridge parameters on the rational four-dimensional core. Which additional nonorthogonal single-bridge parameters occur? A rational Gram-form criterion must handle singular forms as well as positive definite ones. Which whole Schmidt spectra, beyond these bridge spectra and their tensor products, can perfect finite games self-test? The per-ratio existence theorem does not answer that question.

Can any finite game force an irrational ratio? Algebraicity for exclusively infinite-dimensional perfect realizations remains unresolved here. Corollary~\ref{cor:seeds} covers only finite rigid cores with the checked homogeneous generation; Proposition~\ref{prop:totallyreal} concerns a single bridge with its special three-value ideal spectrum. Neither proves general algebraicity or rationality. A general reduction of finite-game type distinctions to universally forced ratios is also absent.

Can a finite game distinguish type $\mathrm{III}_{1/2}$ from type $\mathrm{III}_{2^{-1/2}}$ at the worst-state level? This remains an open question here. Their positive Connes invariants are $2^\Z$ and $2^{\Z/2}$, with the same intersection $2^\Z$ with $\Q_{>0}$. Rational membership tests therefore agree. An odd half-power is needed for a single membership query to distinguish the groups, but is not proved necessary for an arbitrary finite-game distinction.

Determine the sharp robustness constant for $\GIII$. The converse fixes $121.5$ as the natural candidate. No validated numerical constant is given; the sharp constant is conjectural. Even with a sharp robustness law, finite-error remainder control is needed before quoting numerical resource lower bounds. Explicit fixed-state gap constants, in terms of the state and modular spectral distance, are another open problem.

The eight open cells in Section~\ref{sec:phase} include polynomial-price type-III-only games and exponentially expensive games of minimum type $\mathrm I_\infty$. Do type-III-only value-one games admit positive existing exponents below $1/2$, zero exponents, finite exponents above one, or irregular profiles? The realization interval and the limit-computability inclusion do not settle these questions. Determine the exact arithmetical placement of semifinite and type-III-only hosting, and whether a value-one-promise-preserving hardness reduction exists.

\begin{conjecture}[Dilation and rate conjectures]\label{conj:dilation}
For a finite rational type-III-only game with quantum value one, some fixed $r>1$ has $\delta^{\rm sf}_{G,r}(\eps)\to0$ (DC); moreover such a modulus can be chosen with $\delta^{\rm sf}_{G,r}(\eps)\le C\eps^\gamma$ for some $C,\gamma>0$ (RC).
\end{conjecture}
These conjectures would imply an exponential price lower bound by Theorem~\ref{thm:rank-defect}; the proved inequality does not establish them or make them necessary for every exponential-price mechanism. The abstract ladder minimum-type column is the authors' base-game Theorem 13.2 \cite{Base}; its proposed general type-II game-price row remains unproved.

Finally, determine how much noninjective hosts can be distinguished beyond the expected-subfactor criterion. Proposition~\ref{prop:blind} has an injective scope, and the separation above shows that it cannot be extended unchanged. A classification of the entire algebra generated by a perfect strategy is also finer than a classification of ambient hosts.

\subsection{How this work was done}
The authors developed the definitions, proofs and exposition with assistance from AI models from more than one developer, and take responsibility for all claims. The argument combines explicit finite predicates and matrix calculations with named operator-algebraic structure theorems. The robust sector estimate \eqref{eq:sector} is imported from the authors' companion paper; the new gluing and singular-value implementation are presented here. State-orbit, expected-subfactor and ultrapower results are mathematical imports, not conclusions of finite computation.

The ancillary runner executes fourteen checks sequentially on one CPU, enforcing a memory limit below one decimal gigabyte and raising an error on any nonzero exit. Exact rational and quadratic-field computations check predicates, forbidden probabilities, normalization, conference orthogonality, rational forcing chains, all events in the six additional games, nonorthogonal witnesses and graph primitivity. The type/price checks numerically sample the rank/defect estimates and exactly check Boolean tables and tail arithmetic. Independent calculations check the endpoint loss identity, counterexamples to unrestricted profile identities, the sequential-instrument estimate and numerical Airy/Gaussian constants. Floating-point quadrature and optimization are diagnostics, not substitutes for the variational and arbitrary-algebra proofs. Recorded outputs and a SHA-256 manifest accompany the code. They do not formally verify the imported classification results, universal quantifiers, or any physical state-admissibility premise.

\bibliographystyle{abbrv}
\bibliography{references}
\end{document}
